% Français 12357, batch 1 — Gallica views 1–20.
% Pass: Opus 5.5 (claude-opus-5-5), 2026-09-22 — first pass, unchecked against the views by a human.
% Every view was read from the local mirror (archives/fr-12357/), in four
% full-width bands or eight half-width bands at 2400 px, with tight crops for
% numbers and doubtful words, and drafted view by view before the file was
% assembled. No published edition of Mersenne was opened.
%
% What the views are. Greyscale scans, portrait, about 3920 × 5745, one page
% per view. Views 4 and 5 are leaves bound in sideways: they read with the
% view turned a quarter-turn to the left (and the left half of view 5 upside
% down). View 3 is a large folding plate.
%
% Blank views, absent from the file. View 1: the marbled board with the
% printed label « FR. 12357 ». View 2: the flyleaf, blank but for the old
% shelfmark « Suppl.t fr. 902. » and « Volume de 26 Feuillets 3 Août 1894 »
% in a librarian's hand. The gap in the numbering is the record.
%
% What the twenty views hold. Two different things.
%
%   v. 3–5    Three leaves on musical theory in another, older-looking hand
%             than the copy (looped s and e), with no reference to the
%             Harmonie universelle: v. 3, a drawn folding plate — a staircase
%             of 27 string lengths (57600 … 28800), a grid I–XXVII of
%             consonances and dissonances, tables « des consonnes » and « des
%             dissonances » with their ratios, three rows of enharmonic,
%             chromatic and diatonic degrees, and a keyboard of three octaves
%             (« Octaue de vingtsept / vingtcinq / dixneuf marches »); v. 4,
%             a monochord divided in fractions, a solmisation table (B mol,
%             Natur, B dur) and a short text on finding the « feintes » G#,
%             C#, F# by adding and subtracting ratios; v. 5, « Disposition des
%             jntervalles des tons, et semitons appropriés au Clauier selon la
%             gamme jnuentée par Boetius », a keyboard graduated in commas,
%             each interval valued in commas of a 53-comma octave, and a table
%             of concordances and discordances marked + and −.
%   v. 6–7    Music only, on both sides of a leaf: three-part settings headed
%             « Ipolidien », « Dorien », « Frigien » (v. 6), four-part
%             settings with small numbers over the notes (v. 7). Not
%             transcribed; described in notes.
%   v. 8–20   The copy itself: « Remarques tirées du liure de l'harmonie
%             uniuerselle du P. Mersenne ainsy qu'il les auoit escrites de sa
%             main a la marge et aux feuillets blancs deuant et derriere dud.
%             liure » (title at the head of v. 8). One hand throughout.
%             - v. 8–9: remarks written on the blank leaves before the book
%               (« pag. 1re en blanc deuant l'harmonie », « pag. 2me en
%               blanc … »): quadrature and rectification of the circle after
%               Fabre (whose own pages 73, 75, 78–79, 89, 93, 95, 97 are
%               cited), Bouilles and « le Sr Cornu », Guillaume Cambray;
%               gut strings; the square of 3.74586 worked digit by digit;
%               then remarks on « la seconde obseruation qui est apres la
%               table des matiere » (by lines 20–22, 25–26, 38–41).
%             - v. 9–14: « liure 1er de la nature des sons » (pp. 14, 32 …
%               70): speed of sound (1380 pieds, 230 toises a second, the
%               cannon from the Arsenal to the Minimes of Vincennes), weight
%               of air (1700 times lighter than water, against Galileo's
%               400), echoes (the Tuileries, an echo in Descartes's garden),
%               refraction.
%             - v. 11–15: « liure second du mouuement des corps » (pp. 86 …
%               135): fall in the void, Galileo's tables, the path of a body
%               falling to the centre of a turning earth — a spiral, not
%               Galileo's semicircle; the area of the spiral is 8/15 of the
%               sector, with the lemma on sums of squares worked in a
%               marginal table (v. 11–12); a proof in letters b, e, s that
%               the path is no circle (v. 13); Fermat's helix named (v. 12);
%               spirals of any degree (v. 14); the odd-number law explained
%               (v. 14); resistance of air and water after Galileo's
%               « nouuelles pensées » (v. 15); pendulums of various shapes
%               against the isochronous thread (v. 15–16).
%             - v. 16–19: « 3me liure du mouuement » (pp. 157 … 226):
%               vibrating strings and the ellipse, the return of a bent
%               string (v. 17), quadruple tension for the octave (v. 18),
%               weights for the eight degrees of the octave, breaking of
%               cylinders (« ce nombre me semble faux », v. 18), missiles
%               rising and falling, Galileo on the parabolic string.
%             - v. 19: « Du 1er liure de la voix » (pp. 33, 74, 80): the
%               triple progression 1 … 2187; « independant » and infinite
%               perfection; M. de Villiers's dissection of a deaf-mute's ear,
%               1638.
%             - v. 19–20: « Du liure second des chants » (pp. 95 … 111):
%               the Feuillants' psalm-tone (music, « Dixit dominus… »), the
%               number of words from 23 letters after Guldin, the variations
%               of magic squares of 14 and 22, and the rule for listing
%               permutations (DIEU).
%             The last remark of the batch ends at the foot of v. 20; v. 21
%             opens a new remark, « p. 113. au commencemt ».
%
% How the references are written. Each remark opens with the page of the
% Harmonie universelle (« p. 32 », « pag. 38 », « pa. 169 »), then the place
% on that page (« au droit des lignes 23 iusques a 32 », « lig. 22 apres ces
% mots (… », « au droit de la figure », « sur le corollaire premier », « a la
% prop. xxixme »). They are transcribed exactly as written and never checked
% against the book. The page numbers run upwards within each treatise and
% start again with each new one (v. 19: p. 226, then p. 33 of the next
% treatise). The copy names the treatise in the left margin when it changes;
% between those marginal heads the attribution of a page to a treatise is the
% copy's own. The remark on v. 8–9 also cites pages of Fabre's book, which are
% not pages of the Harmonie. Cross-references inside the remarks name other
% parts of the book or other books: « la 9 prop. de l'vtilité [de
% l'harmonie] », « la 1re obseruation physique a la fin de ce volume », « le
% 3me liure de la geometrie de mr Descartes », « Apiarium progymn. », « le
% liure latin de l'harmonie », « le gros volume manuscrit entier de ceste
% matiere ».
%
% The hand of the copy (v. 8–20). One fast, regular seventeenth-century
% cursive, reading end to end at band scale; only numbers and letters at the
% gutter need crops. Kept as written: « co.e » (comme), « p.r » (pour),
% « touss.rs » (tousiours), « -m.t » (-ment), « plus.rs », « f.re » (faire),
% « cæt. », « dud. », « lad. », « Sr », « mr ». A superscript is set with
% \textsuperscript. Seconds of time are written with a double stroke after
% the figure (« 20'' »), set $^{\prime\prime}$. The right-hand edge of the
% versos runs into the binding (v. 11, 13, 15, 17): line ends there are
% hidden, completed in \add{} where the sense makes them certain, and \ill{}
% otherwise. The long s is set s. Diagrams are described in notes with their
% lettering, never drawn.
%
% Notation. In the algebra of v. 13 the lengths are small letters, squares
% written « bb », « ee », and the double vertical stroke « ∥ » is used both
% for equality and for « est à … comme » in a proportion; it is set \parallel
% and never modernised. The same stroke sets the comparisons of the marginal
% table on v. 12. On v. 4–5 the « feinte » (sharp) is written as a letter
% followed by a z-like hook, set $\sharp$ in the text; the square b (natural)
% is a squared h, set \texttt{h}. The signs of the consonances on v. 3 are
% set with the nearest standard symbols, said once in the note there.
%
% Numbers on the leaves. A single series of small, thin ink figures stands at
% the top right of the rectos: « 3 » (v. 6), « 5 » (v. 10), « 6 » (v. 12),
% « 7 » (v. 14), « 8 » (v. 16), « 9 » (v. 18), « 10 » (v. 20), one per leaf,
% none on a verso. Batches 2 and 3 read the same series on to 26 (v. 55),
% the certificate's « Volume de 26 Feuillets » (v. 2), so it behaves as the
% library's foliation; but it is ink, not pencil, and, as for Latin 17859,
% Français 9115 and NAF 5176, no \folio{} is written until a human has
% looked (the coordinator's reconciliation of the three batches). On v. 8, where « 4 » is expected, a « (4 » stands
% at the end of the third line of the title, inside the text block and in a
% heavier stroke: it is recorded in a note and not written as \folio{}. No
% figure was found on v. 3, 4 or 5 (the corners of v. 4–5 are covered by the
% next leaf or turned). Other numbers: « 331 » in ink above the title of v. 8;
% the old shelfmark « Suppl.t fr. 902. » on v. 2; the round stamp
% « BIBLIOTHEQUE IMPERIALE MSS » on v. 5 and v. 8. The copy has no
% pagination of its own. Every « p. », « pag. », « pa. » inside the text is a
% page of the printed book (or, on v. 8–9, of Fabre's), never a page of this
% manuscript.

\documentclass[11pt,a4paper]{article}
% The preamble shared by every transcription of the mathematicians' manuscripts in Gallica.
%
% It rests on one idea: **uncertainty compounds, it does not resolve.** A
% transcription that smooths over a doubtful word has destroyed the only thing
% it was made for — a reader must be able to tell, at every point, what was on
% the page and what was supplied. The apparatus macros below are therefore the
% critical apparatus, not decoration.
%
% The same commands are recognised by `scripts/render.mjs`, which produces the
% site's reading view, and by `scripts/tei.mjs`. A macro added here must be
% added there too, or rendering fails loudly — which is the wanted behaviour.
%
% Comments here are English, like the rest of the repository. The documents
% this preamble sets are French, and that is deliberate: see the skills.

\ProvidesPackage{germain}[2026/09/15 Transcriptions of mathematicians' manuscripts in Gallica]

\usepackage{amsmath,amssymb,amsthm}
\usepackage{mathrsfs}
% Kept from the parent preamble so the renderer's diagram support stays
% available; these manuscripts draw no commutative diagrams, and nothing here uses it.
\usepackage{tikz-cd}
\usepackage[english,french]{babel}
\usepackage{fontspec}

% --- Greek ------------------------------------------------------------------
% The text face stays fontspec's default, Latin Modern, which has no Greek:
% XeTeX drops every glyph it lacks with a line in the log and nothing on the
% page, and the Greek words the manuscripts quote vanished from the PDFs. So
% the Greek and Greek Extended blocks get a XeTeX character class of their own,
% and entering or leaving a run of it switches the family to CMU Serif — the
% same Computer Modern design, with polytonic Greek — and back. Only the
% family changes: a Greek word in \textit or \textbf keeps its shape and series.
% The transitions to and from every other class carry the tokens that class
% has towards an ordinary letter, so French spacing before « ; » or inside
% « » still holds after a Greek word. No grouping, so no brace or \endgroup
% can come out unbalanced; maths is left alone, since XeTeX inserts nothing
% there. Kept identical in `exercices.sty`.
\makeatletter
\newfontfamily\germainGreekfont{cmun}[
  Extension=.otf, NFSSFamily=germaingreek,
  UprightFont=*rm, ItalicFont=*ti, SlantedFont=*sl,
  BoldFont=*bx, BoldItalicFont=*bi, BoldSlantedFont=*bl]
\newXeTeXintercharclass\germain@greek
\def\germain@greekrange#1#2{%
  \@tempcnta=#1\relax
  \loop
    \XeTeXcharclass\@tempcnta=\germain@greek
    \ifnum\@tempcnta<#2\advance\@tempcnta\@ne\repeat}
\germain@greekrange{"0370}{"03FF}
\germain@greekrange{"1F00}{"1FFF}
\def\germain@greekprev{\rmdefault}
\def\germain@greekin{%
  \edef\germain@greekprev{\f@family}\fontfamily{germaingreek}\selectfont}
\def\germain@greekout{\fontfamily{\germain@greekprev}\selectfont}
\def\germain@greekjoin{%
  \XeTeXinterchartokenstate=\@ne
  \@tempcnta=\z@
  \loop
    \ifnum\@tempcnta=\germain@greek\else
      \XeTeXinterchartoks\germain@greek\@tempcnta=\expandafter{%
        \expandafter\germain@greekout\the\XeTeXinterchartoks\z@\@tempcnta}%
      \XeTeXinterchartoks\@tempcnta\germain@greek=\expandafter{%
        \the\XeTeXinterchartoks\@tempcnta\z@\germain@greekin}%
    \fi
    \ifnum\@tempcnta<4095\advance\@tempcnta\@ne\repeat}
% babel-french sets its own classes, and saves and restores the interchar
% state, each time a language is selected: join again after it does.
\addto\extrasfrench{\germain@greekjoin}
\addto\extrasenglish{\germain@greekjoin}
\AtBeginDocument{\germain@greekjoin}
\makeatother

\usepackage{geometry}
\geometry{a4paper,margin=25mm,marginparwidth=28mm}
\usepackage{marginnote}
\usepackage{xcolor}
\usepackage[normalem]{ulem}
\setlength{\emergencystretch}{3em}
\usepackage{hyperref}
\hypersetup{colorlinks=true,linkcolor=black,urlcolor=black,pdfborder={0 0 0}}

% --- Batch metadata ---------------------------------------------------------
% Taken as they stand by the rendering script, which shows them at the head of
% the reading view. They fix what the file claims to cover. The macro names are
% the parent project's, so that the scripts stay shared; \folder here means the
% volume's slug — `fr-9115` — and \pages counts Gallica views.

\newcommand{\folder}[1]{\gdef\grothFolder{#1}}
\newcommand{\batch}[1]{\gdef\grothBatch{#1}}
\newcommand{\pages}[2]{\gdef\grothFirst{#1}\gdef\grothLast{#2}}
\newcommand{\foldertitle}[1]{\gdef\grothFoldertitle{#1}}

% The holder's own shelfmark, printed as they write it: « Français 9115 ».
\newcommand{\shelfmark}[1]{\gdef\grothShelfmark{#1}}

% The dating, copied verbatim from the holder's catalogue — never inferred
% here. « XIXe siècle » is the BnF's; a pass that dates a leaf from its content
% says so in a \note{}, as its own inference.
\newcommand{\dating}[1]{\gdef\grothDating{#1}}

% The status notice, printed in the title block and repeated as a diagonal
% watermark on every page of the PDF. The manuscripts are in the public
% domain; what this line declares is not a right but a quality — an
% unverified first pass by a machine. The skills write the line; a file
% without it gets no watermark, so leaving it out is visible in the output.
\newcommand{\watermark}[1]{\gdef\grothWatermark{#1}}

\gdef\grothFolder{?}\gdef\grothBatch{}\gdef\grothFirst{?}\gdef\grothLast{?}%
\gdef\grothFoldertitle{}\gdef\grothShelfmark{}\gdef\grothDating{}\gdef\grothWatermark{}

% --- The résumé -------------------------------------------------------------
\newenvironment{resume}
  {\small\setlength{\parskip}{0.45em}}
  {\par\medskip\hrule\medskip}

% --- Keywords ---------------------------------------------------------------
% In English, closing the modernised reading's résumé: the modern vocabulary
% under which to search for what the leaves construct. The manifest extracts
% them to tag the volume — this line is the single source of the tags.
\newcommand{\keywords}[1]{\par\smallskip\noindent
  {\footnotesize\textsc{Keywords} --- \textit{#1}}\par}

% --- The critical apparatus -------------------------------------------------

% The Gallica view number, in the margin. This is the anchor that lets a
% disagreement over a formula be settled by looking at *one* image rather than
% a volume of 750. The site uses it to turn the facsimile as the transcription
% is scrolled. A view is one scanned image: usually one side of a leaf.
\newcommand{\page}[1]{%
  \marginnote{\footnotesize\color{gray!70}\textsf{v.\,#1}}%
  \label{page:#1}\ignorespaces}

% The BnF's pencilled foliation, where the leaf carries it and a pass can read
% it: « 198r ». This is what the literature cites — « ff. 198r–208v » — and
% the one line that turns a citation into a view. Recorded, never inferred:
% a folio a pass cannot read is not written.
\newcommand{\folio}[1]{%
  \marginnote{\footnotesize\color{gray!50}\textsf{f.\,#1}}\ignorespaces}

% The modernised reading's coarser counterpart to \page{}: which views a
% section is drawn from.
\newcommand{\pagerange}[2]{%
  \marginnote{\footnotesize\color{gray!70}\textsf{v.\,#1--#2}}%
  \ignorespaces}

% Illegible: never guessed.
\newcommand{\ill}{\textcolor{red!60!black}{[\,\ldots\,]}}

% A doubtful reading: the word is given, and flagged as uncertain.
\newcommand{\uncertain}[1]{\uline{#1}}

% An editorial addition — a missing letter, an implied word.
\newcommand{\add}[1]{\textcolor{blue!50!black}{[#1]}}

% Struck out by the writer. What was crossed out is part of the document.
\newcommand{\struck}[1]{\sout{#1}}

% The transcriber's note, on the state of the page or the mathematics.
\newcommand{\note}[1]{\par\smallskip{\small\color{gray!60!black}\textit{[#1]}}\par\smallskip}

% A marginal note of hers, distinct from the body of the text.
\newcommand{\marginal}[1]{\par\smallskip\noindent\hspace{1em}%
  {\small\color{gray!60!black}#1}\par\smallskip}

% --- Title ------------------------------------------------------------------
% The title block is French, like the documents themselves.

\AddToHook{shipout/background}{%
  \ifx\grothWatermark\empty\else
    \begin{tikzpicture}[remember picture, overlay]
      \node[rotate=52, scale=3.4, align=center, text=gray!18,
            font=\sffamily\bfseries]
        at (current page.center) {\grothWatermark};
    \end{tikzpicture}%
  \fi}

\AtBeginDocument{%
  \begin{center}
    {\large\bfseries Manuscrits de mathématiciens (Gallica) ---
      \ifx\grothShelfmark\empty\grothFolder\else\grothShelfmark\fi}\\[2pt]
    {\itshape\grothFoldertitle}\\[4pt]
    {\small vues \grothFirst--\grothLast
      \ifx\grothBatch\empty\else\ (lot \grothBatch)\fi}\\[2pt]
    \ifx\grothDating\empty\else
      {\small\color{gray!60!black}Datation du catalogue : \grothDating}\\[2pt]
    \fi
    \ifx\grothWatermark\empty\else
      {\small\bfseries\color{red!45!gray}\grothWatermark}\\[2pt]
    \fi
    {\footnotesize\color{gray!60!black}Transcription établie d'après les images IIIF de Gallica.
     Source gallica.bnf.fr / Bibliothèque nationale de France.\\
     Les lectures incertaines sont soulignées, les illisibles notées [\,\ldots\,],
     les ajouts entre crochets ; \textsf{v.} numérote les vues de Gallica,
     \textsf{f.} la foliotation au crayon de la BnF.}
  \end{center}
  \vspace{1em}}

\endinput


\folder{fr-12357}
\shelfmark{Français 12357}
\batch{1}
\pages{1}{20}
\dating{1601-1700}
\watermark{Lecture automatique\\non vérifiée}
\foldertitle{« Remarques tirées du livre de l'Harmonie universelle du P. MERSENNE, ainsy qu'il les avoit escrites, de sa main, à la marge et aux feuillets blancs devant et derrière dudit livre ».}

\begin{document}
\page{3}
\note{Grande planche dépliante, dessinée à la plume et non gravée ; elle porte les traces de ses plis. Elle réunit : en haut, une échelle de vingt-sept cordes en escalier, chacune avec sa longueur ; au milieu, une grille carrée dont les lignes et les colonnes sont numérotées en chiffres romains de I à XXVII, chaque case portant l'un des signes des deux tables de gauche (consonances et dissonances) ou un chiffre ; sous la grille, trois rangées « Degrés Enharmoniques », « Degrés Chromatiques », « Degrés Diatoniques » qui placent les lettres des notes sous les colonnes ; à droite, une colonne verticale de touches marquées des lettres C D E F G A \texttt{h} et de leurs variantes (d, e, f, g, a, b) ; en bas, un clavier en perspective de trois octaves. La grille n'est pas transcrite case par case ; on donne ci-dessous tout le texte de la planche. Un petit trait oblique en haut à droite de la planche n'est pas un chiffre lisible.}
\page{3}
\note{Grande planche dépliante, dessinée à la plume et non gravée ; elle porte les traces de ses plis. Elle réunit : en haut, une échelle de vingt-sept cordes en escalier, chacune avec sa longueur ; au milieu, une grille carrée dont les lignes et les colonnes sont numérotées en chiffres romains de I à XXVII, chaque case portant l'un des signes des deux tables de gauche (consonances et dissonances) ou un chiffre ; sous la grille, trois rangées « Degrés Enharmoniques », « Degrés Chromatiques », « Degrés Diatoniques » qui placent les lettres des notes sous les colonnes ; à droite, une colonne verticale de touches marquées des lettres C D E F G A \texttt{h} et de leurs variantes (d, e, f, g, a, b) ; en bas, un clavier en perspective de trois octaves. La grille n'est pas transcrite case par case ; on donne ci-dessous tout le texte de la planche. Un petit trait oblique en haut à droite de la planche n'est pas un chiffre lisible.}

A \quad La longueur AB, estant doublée donnera la plus grande corde du Sisteme \quad B
\note{Écrit verticalement le long de la première bande, à gauche de l'escalier des cordes.}

Longueurs des vingt-sept cordes, de la plus longue à la plus courte, telles qu'elles sont écrites (verticalement) dans l'escalier :
57600, 55296, 54000, 51840, 51200, 50625, 50000, 48600, 48000, 46080, 45000, 43500, 43472, 40960, 40500, 40000, 38400, 36864, 36000, 34560, 33750, 32400, 32000, 30720, 30375, 30000, 28800.
\note{Le titre de cette ligne est de l'éditeur ; la planche ne donne que les nombres. « 43500 » et « 43472 », les douzième et treizième, sont lus tels qu'ils sont tracés et vérifiés à fort grossissement.}

\subsection*{Table des consonnes}
\begin{itemize}
  \item[$\ast$] Tierce mineure \quad $\frac{5}{6}$
  \item[$\dagger$] Tierce maieure \quad $\frac{4}{5}$
  \item[$\square$] Quarte \quad $\frac{3}{4}$
  \item[$\triangle$] Quinte \quad $\frac{2}{3}$
  \item[$\#$] Sexte mineure \quad $\frac{5}{8}$
  \item[$\star$] Sexte maieure \quad $\frac{3}{5}$
  \item[$\infty$] Octaue \quad $\frac{1}{2}$
\end{itemize}
\note{Chaque consonance a son signe devant elle, celui qui sert dans la grille : une étoile à six branches (tierce mineure), une croix en forme de t barré (tierce majeure), un petit carré (quarte), un triangle (quinte), une croix à traits multipliés (sexte mineure), une étoile (sexte majeure), une boucle ouverte (octave). Les signes imprimés ici sont les plus proches dont dispose la composition. Les rapports sont écrits en fractions verticales dans une colonne à droite.}

\subsection*{Table des dissonances}
\begin{itemize}
  \item[$\cdot$] Comma mineur \quad $\frac{2025}{2048}$
  \item[o] Comma maieur \quad $\frac{80}{81}$
  \item[1] Diese \quad $\frac{125}{128}$
  \item[--] Demiton sous minime \quad $\frac{243}{250}$
  \item[2] Demiton minime \quad $\frac{625}{648}$
  \item[3] Demiton mineur \quad $\frac{24}{25}$
  \item[4] Demiton moyen \quad $\frac{128}{135}$
  \item[5] Demiton maieur \quad $\frac{15}{16}$
  \item[6] Demiton maxime \quad $\frac{25}{27}$
  \item[7] Ton mineur \quad $\frac{9}{10}$
  \item[8] Ton maieur \quad $\frac{8}{9}$
  \item[9] Fausse quarte \quad $\frac{75}{96}$
  \item[10] Triton \quad $\frac{32}{45}$
  \item[11] Fausse quinte \quad $\frac{45}{64}$
  \item[12] Quinte Superflüe \quad $\frac{48}{75}$
  \item[13] Septiesme mineure \quad $\frac{148}{225}$
  \item[14] Septiesme maieure \quad $\frac{10}{18}$
  \item[15] Semi octaue \quad $\frac{25}{48}$
\end{itemize}
\note{« 148/225 » pour la septième mineure est lu tel qu'il est écrit ; le premier chiffre après 1 est un 4 net. Les signes des dissonances (point, o, tiret, chiffres 1 à 15) sont ceux de la grille.}

\note{Sous la grille, trois rangées titrées « Degrés Enharmoniques », « Degrés Chromatiques », « Degrés Diatoniques ». Au pied du clavier, trois accolades : « Octaue de vingtsept marches », « Octaue de vingtcinq marches », « Octaue de dixneuf marches ». Les touches du clavier portent, sur le rang inférieur, C \texttt{h} A G F E D C \texttt{h} A G F E D C \texttt{h} A G F E D C, de droite à gauche, et sur les rangs supérieurs les lettres des touches feintes (d, e, f, g, a, b et leurs doubles, capitales et minuscules). Le \texttt{h} est ici la lettre de la note entre A et C, tracée comme un h à haste barrée.}

\page{4}
\note{Feuillet monté en travers : l'écriture court du haut vers le bas de la vue, et se lit en tournant la vue d'un quart de tour vers la gauche. Main différente de celle de la copie des remarques (vues 8 et suivantes) : une cursive à l'ancienne, aux s et aux e bouclés. Toute la moitié gauche est une figure : une ligne horizontale, le monocorde, divisée de G à un point marqué « 2 » à son extrémité droite, avec au-dessus les lettres des notes et au-dessous la longueur de chaque note en fraction ; au-dessous de la ligne, une famille d'arcs de cercle issus de G, chacun marqué de la partie aliquote qu'il détermine ; puis une table de solmisation et une échelle des demi-tons. Ce qui suit donne le texte de la figure dans l'ordre de la ligne.}

\note{Sur la ligne, de gauche à droite : G ; $\ast$ $\frac{128}{135}$ ; A $\frac{8}{9}$ ; B $\frac{5}{6}$ ; \texttt{h} $\frac{4}{5}$ ; C $\frac{3}{4}$ ; $\ast$ $\frac{96}{135}$ ; D $\frac{2}{3}$ ; b. $\frac{5}{8}$ ; E $\frac{3}{5}$ ; F $\frac{9}{16}$ ; $\ast$ $\frac{8}{15}$ ; G $\frac{1}{2}$ ; $\ast$ $\frac{64}{135}$ ; A $\frac{4}{9}$ ; b $\frac{5}{12}$ ; \texttt{h} $\frac{2}{5}$ ; $\ast$ $\frac{16}{45}$ ; puis, sans fractions, D, $\ast$ (avec « b. » au-dessus), e, f, $\ast$, G ; et « 2 » au bout de la ligne. Le signe noté ici $\ast$ est la feinte, un sautoir hachuré ; le \texttt{h} est le bécarre, tracé comme un h carré. Sous la ligne, les arcs portent « $\frac{1}{3}$ », « $\frac{1}{3}$ », « $\frac{1}{4}$ », « $\frac{1}{5}$ », « $\frac{1}{6}$ », « $\frac{1}{6}$ ou $\frac{1}{7}$ », « $\frac{1}{8}$ ».}

\note{À gauche, une table en trois colonnes, les lettres G F E D C B A G F E D C B A G en marge, avec pour chaque lettre sa syllabe : première colonne, titrée au pied « B mol » : re, ut, fa, la, sol, fa, mi, re, mi, fa, la, sol, fa, mi, re, ut (plusieurs syllabes récrites l'une sur l'autre dans le haut de la colonne) ; deuxième colonne, « Natur » : sol, fa, mi, re, ut, \texttt{h}, la, sol, fa, mi, re, ut, \texttt{h}, la, sol, fa ; troisième colonne, « \texttt{h} quarre ou B dur » : ut, (vide), la, sol, fa, mi, re, ut, (vide), la, sol, fa, mi, re, ut. À droite de la table, deux échelles descendantes de signes « o $\ast$ o », « b.o$\ast$o », une par degré, et deux colonnes de noms d'intervalles : « Lima semitonium », « Lima semiton », « \uncertain{diese} semiton », « semiton », « Lima semiton », « \uncertain{Diese} semiton », « semiton », « Lima semiton » ; puis « Lima semiton », « \uncertain{diese} semiton », « semiton », « Lima semiton », « \uncertain{diese} semiton », « semiton », « Lima semiton », « Lima semiton ».}

Apres auoir les notes des concordances parfaites au monocorde lesquelles sont denotées par demicercle reste a descrire les autres notes mitoyennes come G$\sharp$. A. C$\sharp$. F. F$\sharp$ pour trouuer A. faut soubtraire la raison harmoniq\add{ue} $\frac{3}{4}$ de D. que fait $\frac{2}{3}$ reste $\frac{8}{9}$ pour A.

Pour trouuer F faut adiouster $\frac{3}{4}$ auec $\frac{3}{4}$ de C qui fait $\frac{9}{16}$ pour F

Pour trouuer la feinte de \uncertain{C} scauoir C$\sharp$ faut adiouster $\frac{4}{5}$ auec A. $\frac{8}{9}$ fait $\frac{32}{45}$ pour C$\sharp$.

Que sy de la C$\sharp$. contenant $\frac{32}{45}$ vous soustraiez $\frac{3}{4}$ la reste sera $\frac{128}{135}$ pour G$\sharp$.

Et sy a C$\sharp$ $\frac{32}{45}$ on adiouste encore $\frac{3}{4}$ la somme sera $\frac{8}{15}$ pour la feinte de F$\sharp$.
\note{Le texte écrit la feinte par la lettre suivie d'un z à queue descendante, attaché à la lettre ; ce signe est rendu $\sharp$. « soubtraire » et « adiouster » sont à entendre au sens des raisons : les fractions se multiplient ou se divisent. Au-dessus de la fin de la première ligne, un petit « 5 » suivi d'un signe, sans rapport visible avec le texte. Le grand demi-cercle de la figure traverse la fin des lignes 5 à 9 ; il ne biffe rien.}

\page{5}
\note{Feuillet de même papier et de même main que celui de la vue 4, monté lui aussi en travers et à demi recouvert, en haut de la vue, par le bord du feuillet suivant (on y voit des portées de musique et le mot « Ipolide », qui sont de la vue 6). Il porte deux choses écrites dans deux sens : dans la moitié droite de la vue, qu'on lit en la tournant d'un quart de tour vers la gauche, un titre et un clavier gradué ; dans la moitié gauche, qu'on lit en retournant la vue tête-bêche, une liste des intervalles et une table. Le cachet rond de la Bibliothèque impériale (« BIBLIOTHEQUE IMPERIALE MSS ») est posé sur le clavier. Ce feuillet et le précédent ne sont pas des remarques renvoyant au livre : aucun renvoi à une page de l'\emph{Harmonie universelle} n'y figure.}

Disposition des jntervalles des tons, et semitons appropriés au Clauier selon la gamme jnuentée par Boetius pour la reformation de la musique soubs Gregoire premier, et depuis continuée en l'Eglise Romaine jusques au present auec \ill{} de chant Gregorien
\note{La fin de la deuxième ligne passe sous le bord du feuillet suivant ; un ou deux mots y sont cachés.}

\note{Sous le titre, un clavier dessiné, surmonté d'une règle graduée de 0 à 100, chiffrée de dix en dix. Chaque touche porte en tête le nombre de commas qu'elle occupe sur la règle ; de gauche à droite : 4, 5, 5, 3, 5, 4, 5, 5, 3, 5, 4, 5, 4, 5, 5, 3, 5, 4, 5, 5, 3, 5, 4, 5 (le cinquième chiffre, « 5 », est récrit sur un autre). Les touches hachurées, les feintes, sont marquées « Fe », « Ge », « b. », « Ce », « b. », « Fe », « Ge », « b. », « Ce », « b », « Fe » ; les touches blanches, en bas, « G A \texttt{h} C D E F G A \texttt{h} C D E F ».}

\note{Ce qui suit est écrit tête-bêche par rapport à ce qui précède. La première ligne passe elle aussi sous le bord du feuillet suivant : on n'en voit que le bas, avec le dénominateur « 81 » d'une fraction, écrit au-dessus de « la diuision » ; le numérateur et les mots qui précèdent sont cachés.}

\ill{} de la longueur du monocorde, en la diuision du present clauier est $\frac{1}{53}$ partye de loctaue

Diese contient les $\frac{24}{25}$ de la longueur du monocorde et dans le present clauier ne uaut que 3. comma

Lima contient $\frac{128}{135}$ du monocorde, et au present clauier uaut --- 4. comma

Semiton est la $\frac{15}{16}$ partye du monocorde et en ce clauier uaut --- 5 comma

Ton Mineur contient les $\frac{9}{10}$ du monocorde et icy est composé de la diese et du semiton qui font --- 8. comma

Ton parfaict contient les $\frac{8}{9}$ du monocorde et est composé d'un lima et d'un semiton faisan\add{t} 9. com\add{ma}

Ton Mayeur ou excedant contient $\frac{225}{256}$ du monocorde et est composé de 2 semitons faisant --- 10 comma

Tierce mineure contient $\frac{5}{6}$ du monocorde et est Composé d'un ton parfait et d'un semiton ou --- 14 com\add{ma}

Tierce Mayeure contient les $\frac{4}{5}$ du monocorde et est composée d'un ton parfait et d'un ton mineur ou de --- 17 comma

Quart contient les $\frac{3}{4}$ du monocorde et comprend un ton parfait un ton mineur et un semiton ou --- 22 comma

Quinte contient les $\frac{2}{3}$ du monocorde et comprend 2 tons parfait un ton mineur et un semiton ou --- 31 co\add{mma}

Sexte mineure contient les $\frac{5}{8}$ du monocorde et comprend 2 tons parfait 1 \uncertain{tons} maieur, et un ton mineur ou --- 36 comma

Sexte maieure contient les $\frac{3}{5}$ du monocorde et comprend 2 tons parfait 2 tons mineur et un semiton ou --- 39. co\add{mma}

Octaue contient la $\frac{1}{2}$ du monocorde et comprend 3. ton parfait 2 tons mineur et un ton maieur ou --- 53 comma
\note{Les tirets rendent un long trait de plume tiré jusqu'au nombre de commas, en fin de ligne. Plusieurs fins de mot (« faisan », « com », « co ») sont perdues au bord déchiré du feuillet ; les compléments sont entre crochets.}

La table suiuante denotte les concordances et discordances lesquelles sont marquées par le signe $+$ plus et $-$ moins
\[
\begin{array}{ll|l|l|l|l|l|l|l}
 & & \frac{1}{3}\ \text{mineure} & \frac{1}{3}\ \text{maie.} & \text{quart} & \text{quinte} & \frac{1}{6}\ \text{min} & \frac{1}{6}\ \text{maieur} & \frac{1}{8} \\ \hline
\text{G} & & \text{B} & \text{h} & \text{C} & \text{D} & \text{De} & \text{E} & \text{G} \\
 & \text{Ge} & \text{h}-1 & \text{C}+1 & \text{Ce} & \text{De}+1 & \text{E}-1 & \text{F}+1 & \text{Ge} \\
\text{A} & & \text{C}-1 & \text{Ce} & \text{D} & \text{E}-1 & \text{F}-1 & \text{Fe} & \text{A} \\
\text{b.} & \text{B} & \text{Ce}-2 & \text{D} & \text{De} & \text{F}-1 & \text{Fe}-2 & \text{G} & \text{B} \\
\text{h} & & \text{D} & \text{De}+2 & \text{E} & \text{Fe} & \text{G} & \text{Ge}-1 & \text{h} \\
\text{C} & & \text{De} & \text{E} & \text{F} & \text{G} & \text{Ge}-1 & \text{A}+1 & \text{C} \\
 & \text{Ce} & \text{E}-1 & \text{F}+1 & \text{Fe} & \text{Ge} & \text{A} & \text{B}+2 & \text{Ce} \\
\text{D} & & \text{F}-1 & \text{Fe} & \text{G} & \text{A} & \text{B} & \text{h} & \text{D} \\
\text{b.} & \text{De} & \text{Fe}-2 & \text{G} & \text{Ge}-1 & \text{B} & \text{h}-2 & \text{C} & \text{De} \\
\text{E} & & \text{G} & \text{Ge}+1 & \text{A}+1 & \text{h} & \text{C} & \text{Ce}+1 & \text{F} \\
\text{F} & & \text{Ge}-1 & \text{A}+1 & \text{B}+1 & \text{C} & \text{Ce}-1 & \text{D}+1 & \text{F} \\
 & \text{Fe} & \text{A} & \text{B}+2 & \text{h} & \text{Ce} & \text{D} & \text{De}+2 & \text{Fe}
\end{array}
\]
\note{Table recopiée case par case ; « h » y rend le bécarre, tracé comme un h carré. Les en-têtes « $\frac{1}{3}$ », « $\frac{1}{6}$ », « $\frac{1}{8}$ » sont écrits ainsi pour tierce, sixte et octave. Dans la ligne « Ce », colonne de la sixte majeure, le « 2 » de « B$+$2 » est récrit sur un autre chiffre. Dans la dernière ligne, colonne de la quinte, la case porte un C suivi d'une boucle descendante, lu « Ce » avec doute. Dans la ligne « E », l'octave est écrite « F », telle quelle. Sous chacune des colonnes, sauf la première, le mot « comma » est écrit verticalement ; sous la première colonne des intervalles, un simple trait.}

\page{6}
\note{En haut à droite, à l'encre, d'une main fine étrangère au texte, le chiffre « 3 » : c'est le premier nombre lisible de la série qui numérote chaque feuillet à son recto (voir l'en-tête). La page est entièrement de musique notée, sans texte : trois systèmes de trois portées, en notation blanche avec barres de mesure, armures et changements de clés ; au-dessus des sections, les noms de modes « Ipolidien », « Dorien », « Frigien », puis de nouveau « Dorien », « Frigien », « Dorien ». Elle n'est pas transcrite. Le quatrième système est tracé et resté vide.}

\page{7}
\note{Verso de la page précédente : musique notée, sans texte, en deux systèmes de quatre portées (clés d'ut et de fa), avec au-dessus de presque chaque note de la partie supérieure, et de celles des deux parties intermédiaires, un petit nombre entre 1 et 20 (« 19 », « 15 », « 12 », « 15 14 », « 10 », « 12 13 13 »…) ; la page ne dit pas ce qu'ils comptent. Le cinquième rang de portées de chaque système est tracé et vide. Non transcrite.}

\page{8}
\note{Début de la copie, dans la main qui la tient jusqu'à la fin du lot. En haut à gauche de la page, au-dessus du titre et de la même encre, le nombre « 331 ». À la fin de la troisième ligne du titre, après « liure », une grande courbe ouverte et un « 4 », de l'encre du texte ; à la place où les feuillets suivants portent leur numéro, mais tracé plus gras que les chiffres de cette série. Ce « 4 » n'est pas pris pour le folio (voir l'en-tête). Cachet rond « BIBLIOTHEQUE IMPERIALE MSS » au pied de la page, sur la figure.}

\marginal{\ill{}ge 1\textsuperscript{re} deuant \add{l'h}armonie au \add{feu}illets blancs}
\note{Note de la marge gauche, en regard du titre, de la main du texte. Le début de chaque ligne passe dans la reliure.}

\section*{Remarques tirées du liure de l'harmonie uniuerselle du P. Mersenne ainsy qu'il les auoit escrites de sa main a la marge et aux feuillets blancs deuant et derriere dud. liure}

La sphære ayant 7 pour son diametre, la circonference de son plus grand cercle sera 22, et la superficie totale 154 car le diametre multiplié par la circonference donne la surface. Le quarré Isoperimetre aud. cercle, a scauoir qui a 7 pour diametre, a p\textsuperscript{r} son costé $5\frac{1}{2}$, et en sa circonference 22 pieds sy le rond a 14 en diametre, son aire est 154, et le quarré isoperimetre n'en contient que 121, de sorte que le \emph{poids} du rond sera 14, et celuy dud. quarré : et neantmoins tant le rond comme le quarré ont 44 pieds de tour selon Bouilles.
\note{« poids » est souligné. Après « celuy dud. quarré » la ligne s'arrête sur deux points, sans le nombre qu'on attend ; rien n'est gratté.}

Pour faire vn rond egal a vn quarré, diuisés le diametre du rond en 8, et donnés 10 a la diagonale du quarré, l'aire du rond aura $50\frac{4}{7}$ et celle du \struck{cercle} quarré 50 seulem\textsuperscript{t}, et le rond sera encore plus grand que le quarré de $\frac{1}{88}$. Par le moyen de 70 et 99 multipliés en soy mesmes l'on a deux nombres en raison double --- $\frac{1}{4900}$ qui n'est pas perceptible\textsuperscript{$\times$}. Le S\textsuperscript{r} Cornu approche plus pres faisant que le costé du quarré inscrit soit 9, et quart du cercle circonscrit 10
\note{Le trait qui précède « $\frac{1}{4900}$ » ouvre la ligne ; on ne peut dire s'il vaut « moins » ou s'il n'est qu'un trait de liaison. Le renvoi en forme de croix double après « perceptible » répond au même signe tracé dans la marge gauche, au coin d'une petite figure : un carré dont le côté est marqué « 70 » et la diagonale « 99 », avec un arc de cercle.}

Les cordes de boyau sont meilleures et plus fortes quand le mouton ou le boucq est tué en pleyne lune, et toute la graisse et la chair en doibt estre ostée, le pli leur fait perdre leur rondeur, et les tignes et vermisseaux elles souuent \uncertain{cas}, et se rend fausse. \emph{Fabre} se persuade, qu'vne chorde de loup rend celles de mouton presentes fausses, ou leur nuit, parce que les cochons meurent au ventre de la truye sy l'on chastre le verrat geniteur, ou quand il est tué.
\note{« Fabre » est souligné ici et au paragraphe suivant. « cas, » finit la ligne au bord de la page ; le mot n'est pas achevé.}

\marginal{\ill{}ta.}
\note{Dans la marge gauche, entre ce paragraphe et le précédent, un mot dont le début passe dans la reliure. Plus bas dans la même marge, une esquisse très pâle : un triangle, une perpendiculaire et des arcs, sans lettres.}

\emph{Fabre} p. 73 pretend que la ligne AG tirée de l'angle droit sur le mitan de BC, ou AH, est esgale au quart de cercle AB ; ou il arriue aussy que GH est esgale au demidiametre AE, et le S\textsuperscript{r} Cornu pretend que le quart de cercle AB est moindre d'vne 9\textsuperscript{me} partye que AB quart de circonference ; c'est a dire que le costé du quarré AB estant 9, le quart de la circonference AB est 10 : ce qui approche de la verité et l'une et l'autre pratique est assés exacte p\textsuperscript{r} les mechaniques. Le S\textsuperscript{r} Cornu fait le quart de la circonference trop grand d'vn peu, car le polygone circonscript de 97 costés est plus petit que son cercle supposé. Il cite Denis \emph{\uncertain{auis}} conseiller du Roy au balliage de Paris qui trouuoit que la circonference contenoit le diametre trois fois et $\frac{112}{113}$ d'vne septiesme partie dud. diametre. L'aire du cercle vient en plus\textsuperscript{rs} façons 1\textsuperscript{o} en multipliãt le diametre par soymesmes et puis par 11, et puis les diuiser par 14. 2\textsuperscript{o}. la moitie du diametre par la circonference.
\note{« p. 73 » renvoie, d'après la phrase, à une page du livre de Fabre, non à l'\emph{Harmonie universelle}. Le mot souligné après « Denis » peut se lire « auis » ou « anis ». Figure à droite du paragraphe : un cercle circonscrit au carré ABCD (A et B en haut, D et C en bas), E au centre, G sur BC, H sur DC, les droites AG et AH tirées de A, et la diagonale AC en pointillé.}

Page 75. il cite de M\textsuperscript{r} Guillaume Cambray, chancelier de Bourges (qu'il dit auoir fait beaucoup de beaux ouurages) la maniere d'assembler deux cercles ensemble, c'est à dire d'en faire vn esgal aux deux autres : qui est en prenant le demidiametre de A et le portant du centre C en E, et puis sur la perpendiculaire CG porter le demy diametre de B, de C en D, et ioindre DE, par la ligne DE, et l'ouuerture de laquelle a scauoir CG esgale a ED descriuant le cercle CGH, il sera esgal aux deux cercles A et B, comme seroit le quarré de DE esgal aux quarrés de DC et CE, ce qu'estant cognu aux enfans par les elemens d'Euclide, Je voys que Fabre n'estoit ny geometre, ny philosophe. p\textsuperscript{r} la soubstraction d'vn cercle d'vn autre cercle p. 78 et 79 il est fort obscur. Il est bien plus aysé d'oster vn cercle moindre d'vn plus grand, par exemple A de GH en transportant le Rayon de A de C en E comme deuant, et puis le Rayon CH du point E sur le rayon CG, ou il pourra arriuer, comme icy, en D, car il restera CD p\textsuperscript{r} le rayon du cercle, dont A sera osté ; c'est a dire que l'assiette GH$\alpha\beta\gamma\delta$ sera esgale au cercle A, c'est a dire l'anneau ou armille sera esgale au rond.
\note{« Page 75 » et « p. 78 et 79 » renvoient de même au livre de Fabre. Figure à droite, en bas : deux cercles B et A côte à côte, chacun avec un diamètre horizontal ; à droite, deux cercles concentriques de centre C, le plus grand passant par G en haut et H à droite, le plus petit par D en haut, avec E sur le rayon CH et la droite DE ; les lettres $\alpha$, $\beta$ sous les cercles concentriques et $\gamma$, $\delta$ entre le cercle A et eux. Le cachet de la bibliothèque couvre le cercle A. Au pied de la page, à droite, un trait vertical isolé.}

\page{9}
\note{Verso du feuillet précédent. La page est partagée par des traits horizontaux en quatre remarques, chacune précédée dans la marge gauche du lieu du livre d'où elle est tirée. Aucun nombre dans les coins.}

\marginal{pag. 1\textsuperscript{re} en blanc deuant l'harmonie}

pag. 89 on aura vne armille d'vne concauité de telle grandeur qu'on voudra esgale à vn cercle donné, si par exemple du cercle precedent GH$\beta$ l'on veut auoir vn anneau esgal, qui ayt sa circonference \struck{ex} interieure, ou son vuide esgal a lad. circonference GH$\beta$. soit fait ABCD punctué esgal à HG$\beta$, et que l'ouuerture du compas CB punctuée soit poussée du centre M à L, l'anneau CKDNAO et cæt. sera esgal au cercle GH$\beta$, ou BADC ponctué. De mesme l'on reduict l'anneau en cercle, en prenant le rayon de l'anneau ML, et le transportant du point C concaue de l'anneau iusques sur le rayon ML, au point B, qui se trouue la par hazard se rencontrer, et le cercle descrit du rayon MB sera esgal aud. anneau.
\note{« pag. 89 », et plus bas « pag. 95 », « pag. 93 », « pag. 97 », continuent les renvois aux pages du livre de Fabre de la page précédente ; la remarque, d'après la marge, était écrite par Mersenne sur la première page blanche avant l'\emph{Harmonie}. Figure à droite : deux cercles concentriques en pointillé de centre M, deux diamètres perpendiculaires LN et OC (O, A, M, C sur l'horizontale ; L, B, M, D, N sur la verticale), et le carré inscrit BADC en pointillé. Le cercle GH$\beta$ est celui de la figure de la page précédente.}

pag. 95 il marque encore que le cercle inscrit dans le quarré est comme 11 à 14 et le quarré inscrit dans le cercle cor\textsuperscript{e} 7 à 11./
\note{« cor\textsuperscript{e} » : « comme », abrégé.}

\marginal{$\times$ qui est la moitié de BC.}

pag. 93. on reduict vn quarré AD en triangle isopleure, en tirant vn quart de cercle du point A. et rayon AB, a scauoir BEC, et prenant la distance de C \struck{en} à E\textsuperscript{$\times$} en ligne droicte, et la transportant de D en F, \struck{\ill{}} car FD est la moitie du costé du triangle cherché partant DG en est la hauteur, car le double de FD, c'est a dire FH estant tiré, la mesme ouuerture du compas donne les deux autres costés
\note{« à » est écrit au-dessus de « en » biffé ; le renvoi en forme de croix au-dessus de E répond à la note marginale « qui est la moitié de BC ». Un mot court et raturé avant « car » est illisible sous la rature. Figure à droite : le carré ABDC (A et B en haut, C et D en bas), l'arc BEC de centre A, tiré de B à C, F sur CD, le triangle de sommet G au-dessus de B, avec les droites GC, GF et GD.}

pag. 97 on a le contenu du triangle en multipliant le costé en soy, et le prouenu par 13, et le tout diuisé par 30. par exemple si le costé est 8, l'aire sera $27\frac{11}{15}$. On aura la ligne GD, si on multiplie le costé 8 par 13, et le prouenant diuisé par 15 donnera quotient $6\frac{14}{15}$, ou si GD est conu par luy on scaura le costé GF en multipliant GD par 15, et le produit diuisé par 13./

\marginal{pag. 2\textsuperscript{me} en blanc deuant l'harmonie}

\subsection*{Quarrer vn nombre.}

\[
\begin{array}{lrl}
 & 3.\,74586 & \\ \hline
\text{Quarré de 3} \ldots & 9\;4916256436 & \text{Quarrés des simples figures} \\
\text{Rectangle de 3 en 7 double} & 4\;256408096 & \text{rectangles doubles de chasque figure auec sa precedente} \\
\text{Rectangle de 3 en 4 double} \ldots & 2\;4706460 & \text{rectangles doubles auec la seconde precedente} \\
\text{Rectangle de 3 en 5 double} \ldots & 3\;11248 & \text{rectangles doubles auec la troisiesme precedente} \\
\text{Rectangle de 3 en 8 double} \ldots & 4884 & \text{rectangle double auec la quatriesme precedente qui est de 6.} \\
\text{Rectangle de 3 en 6 double} \ldots & 36 & \\ \hline
 & 5563071396 & \\ \hline
 & 140314671396 &
\end{array}
\]
\note{Carré de 3,74586 posé chiffre par chiffre : chaque ligne est écrite en escalier sous les chiffres du nombre, et les chiffres sont donnés ici tels qu'ils se suivent sur la page. Un trait courbe descend en diagonale et sépare, dans les cinq premières lignes, le premier chiffre (9, 4, 2, 3, 4) du reste ; la dernière ligne, « 36 », est à gauche de ce trait. Un « 1 » isolé, sans doute une retenue, est écrit à droite au-dessous de la ligne « 4884 ». Le total, $140314671396$, est bien le carré de $374586$.}

\marginal{en la seconde obseruation qui est apres la table des matiere}

Au droit des \struck{la} lignes 20, 21, 22.
Il vaut mieux dire que le premier mouuem\textsuperscript{t}. de la chorde plus aigue est desunj, et qu'elle a la moitie de ses coups desunis et autant d'vnis
Audroit des lignes 25 et 26
Et que la quarte a trois des coups du son aigu desunis
Au droit des lignes 38, 39, 40, et 41.
Disposés tellem\textsuperscript{t}. que la quarte soit dessous, elles ne s'vnissent qu'au 6\textsuperscript{me} \struck{bat} ou battem\textsuperscript{t}. aigu, et si la quinte est au dessous, au 4\textsuperscript{me} battem\textsuperscript{t}.
\note{Ici le renvoi n'est plus à une page mais aux lignes d'une page, celle de la « seconde observation » placée après la table des matières.}

\marginal{au liure 1\textsuperscript{er} de la nature des sons.}

p. 14 au droit des \struck{la} lignes. 25 iusques à 34.
Toutes les experiences qui sont icy proposées ayant esté faites, monstrent que le son va touciours d'vne esgale vistesse en tout temps et en tout lieu, et qu'il fait 2\ill{} toises dans vne seconde minute, cor\textsuperscript{e} il sera remarqué dans le 3\textsuperscript{me} liure suiuant, et souuent ailleurs, \struck{voicy} \add{voyez} la 6\textsuperscript{me} prop. de l'vtilité, et la 3\textsuperscript{me} obseruation afin de reformer ceste proposition et ses semblables.
\note{« Livre premier de la nature des sons », p. 14, lignes 25 à 34 : renvoi au livre. Après « fait 2 », la suite du nombre entre dans la reliure ; on lit au bord « ota », peut-être « Nota » écrit dans la gouttière. Le mot biffé avant « la 6\textsuperscript{me} » est surchargé et se lit mal ; « voyez » est écrit au-dessus. « cor\textsuperscript{e} » : « comme ».}

\note{Au pied de la page, à droite, isolé : « pag. 32 », suivi d'une lettre coupée par la reliure — la réclame de la remarque qui ouvre la page suivante.}

\page{10}
\note{En haut à droite, « 5 ». Même disposition jusqu'à la fin du lot : chaque remarque commence par la page du livre (« p. 32 », « pag. 38 »…) et souvent par la ligne ou la figure en regard de laquelle Mersenne l'avait écrite, puis donne la remarque ; un petit trait horizontal dans la marge gauche sépare les remarques. Le premier mot de la dernière ligne de chaque remarque est souvent souligné de ce trait. Les pages se suivent en ordre croissant : on est toujours, semble-t-il, dans le premier livre « de la nature des sons » nommé à la page précédente.}

p. 32 au droit des lignes 23 iusques a 32.
voyez la 1\textsuperscript{re} obseruation physique a la fin de ce volume ou ie donne vne autre maniere de peser l'air, et le pese en effect ; il est 1700 fois plus leger que l'eau, et l'obseruation de galilée, qui trouue que l'air est 400 fois plus leger que l'eau, pag. 81 de ses dialogues, ou iournées. voyés n\textsuperscript{re} 15 et 16 articles qui expliquent ses nouuelles pensées
\note{« pag. 81 de ses dialogues » renvoie à Galilée, non à l'\emph{Harmonie}. « n\textsuperscript{re} » est lu tel quel (peut-être « nostre »).}

pag. 38 au droit des lignes 28 iusques à 39.
l'on a la vraye solution de tout cecy dans la 9 propos. de l'vtilité de l'harmonie, ou la veritable vistesse du son est expliquée, lequel fort ou foible va touciours d'esgale vistesse, et fait 1380 pieds dans vne seconde minute, ou 230 toises le bruit du canon depuis l'allée de l'arsenac iusques au Conuent des minimes de vincennes employe 20\textsuperscript{$\prime\prime$} ou $\frac{1}{3}$ de minute, qu'on voye s'il y a 4600 toises

pag. 41 au droit des lignes 6, 7, 8 et 9
voyes le liure entier des chants ou toutes sortes de Combinations sont expliquées, auec des exemples

pag. 43. apres la ligne 21.
voyes la propos. 35 et 36 du premier liure des consonances
et au droit des lignes 41, 42, et 43.
voyes la derniere propos. du liure de l'orgue, ou est la demonstration de la duplication du cube, et le 3\textsuperscript{me}. liure de la geometrie de m\textsuperscript{r}. descartes.
\note{Une accolade dans la marge réunit les quatre premières lignes de cette remarque. Une tache d'encre après « 36 ».}

pag. 44 \struck{\ill{}} au droit des lignes 35 iusques a 40
voyes la dioptrique et les meteores, ou il est dit que la lumiere n'est rien que le mouuem\textsuperscript{t} de la matiere subtile en droite ligne, lequel mouuem\textsuperscript{t} affectant les bouts du nerf optique s'apelle lumiere.

p. 46. ligne 36. apres ces paroles | parallelos du soleil
sur quoy voyez la marge de la 28 page de l'vtilité de l'harmonie.

p. 50 ligne 31 apres ces mots | a l'octaue plus haute.
J'en ay remarqué vn au haut de nos 2 et 3 degré du 2\textsuperscript{me}. estage qui respond vne octaue plus haut que la voix du P. Mersenne plus creuse (c'est a dire aux minimes) et mons.\textsuperscript{r} des Cartes en a ouy vn en son iardin qui respondoit à toute sorte de bruits comme le cri d'vn poullet, ce qui est estrange, parce que le son plus aigu ne se fait que par vn battement ou mouuem\textsuperscript{t}. d'air plus viste. Refaisant l'essay ie trouuay que ma voix plus haute d'vne quarte que ma plus basse, il respond \quad le 1\textsuperscript{r}, 2 et mieux sur le 3\textsuperscript{me} degré à la 12\textsuperscript{me} ou 5\textsuperscript{me}.
\note{« vn » : un écho, d'après la suite. Un blanc d'environ un tiers de ligne est laissé après « il respond », sans rature : la copie a sans doute laissé la place d'un mot qu'elle n'a pas lu. « voix du P. Mersenne » est bien dans la copie, à la troisième personne.}

p. 58 sur le corollaire premier
voyes la vistesse du son dans la 9\textsuperscript{me} propos. de l'vtilité, et corriges tout cecy, et ce qui est dans le 1\textsuperscript{er} corollaire de la 21\textsuperscript{me} propos. du 3 l. et ailleurs selon lad. 9\textsuperscript{me} propos.
\note{« du 3 l. » : du troisième livre.}

p. 59. sur la prop. 28\textsuperscript{me}
l'echo des tuilleries est rond et non pas en ouale

p. 60. au droit de la figure, de la parabole.
Au point du foyer e, il n'est pas necessaire que plus\textsuperscript{rs} partyes d'air sy assemblent, il suffit qu'il se meuue plus viste cor\textsuperscript{e} font les partyes subtiles p\textsuperscript{r} faire le feu.

p. 61. au droit d'vne figure de parabole
voyez la 5 propos. du liure de la verité, ou toutes ces figures sont faites a la perfection

p. 63. a la prop. xxix\textsuperscript{me}
voyes la lettre de m\textsuperscript{r}. deschamps sur ceste refraction

p. 65. au droit de la figure.
Ce traitté de dioptrique est imprimé voyes la 2\textsuperscript{me} et 3\textsuperscript{me} proposition de celle d'herigonne

p. 67. au droit de la ligne 17.
voyes la premiere obseruation ou il est dit que les cercles de l'air sont 1380 plus vistes que ceux de l'eau/

\page{11}
\note{Verso du feuillet 5. Aucun nombre dans les coins. Le bord droit de la page passe dans la reliure : la fin de plusieurs lignes y est cachée, et les compléments que le sens rend sûrs sont donnés entre crochets ; les autres restent en \emph{lacune}. Contre le bord droit de la vue se voit la table de nombres de la marge gauche de la page suivante ; elle est transcrite à la vue 12.}

p. 69. au commencem\textsuperscript{t}
puis que l'air cede 1380 fois plus aysem\textsuperscript{t} que l'eau, a scauoir sy le co\add{up} d'harquebuse, ou autre coup d'vn corps plus pesant que l'eau ira moin\add{s} loing \struck{\ill{}} 1380 fois dans l'eau que dans l'air
Et sur la fin de lad. page / experience à faire
\note{Le mot biffé après « loing » est couvert d'une tache d'encre.}

p. 70. Corollaire 3\textsuperscript{me}. / experience à faire

\marginal{au liure second du mouuem\textsuperscript{t}. des corps}

p. 86. lig. 7\textsuperscript{me}.
Il faut supposer que la cheute des corps pesans se face dans le vuyde o\add{u} dans vn milieu qui n'empesche nullem\textsuperscript{t}, et cela posé tout corps desce\ill{} aussy viste l'vn que l'autre : or Galilée suppose cela dans ses derniers dialogues

p. 87. lig. 41 apres ces mots (et la 4 D en brasses)
De sorte que ces deux dernieres colonnes contiennent les cheutes suyuant la supposition de Galilée, et les deux premieres suyuant les experiences du pere Mersenne.

p. 93. au droit de la ligne \struck{\ill{}\textsuperscript{me}} 11\textsuperscript{me}
Galilée a \uncertain{dy depuis} auoüé par lettre expresse qu'il auoit manqué en ceste ligne \emph{force} par la demonstration de la marge aux pages 94 et 95, la page 96 et 97 contient encores vne autre demonstration de la mesme chose
\note{Le premier nombre de la ligne est biffé et couvert d'une tache ; « 11\textsuperscript{me} » est écrit au-dessus. « force » est souligné : c'est le mot du livre en cause. Une autre tache d'encre sur « depuis ».}

p. 94 au commencement
\note{Figure sur toute la largeur de la page : un quart de cercle ABKC de centre E sur le diamètre horizontal LEA (L à gauche, A à droite), avec un second arc concentrique très voisin du premier du côté de A ; des rayons en pointillé issus de E aboutissent aux points B, F, G, H, I, K de l'arc et, au-delà de K, à des points de l'arc gradués 27, 28, 29, 30, 40, 41, 42, 43, 44, puis C. Une spirale en pointillé part de A et passe par D, 26, 25, 24, 23, 21, 20 (?), 18, 15, 12, 9, 6, 3 pour finir en E ; de petits arcs de cercle concentriques, en trait plein, relient des points numérotés de 1 à 22 et des lettres L, M, N, O, P, Q, R, S, T. Plusieurs chiffres sont récrits (« 24 », « 25 », « 26 » près de H, G, F).}

Soit ABKC l'equateur dans la terre, et qu'vn corps pesant commence a descendre du poi\add{nt} \ill{} iusques au centre E, dans vn temps qui soit au temps de 24 heures, cor\textsuperscript{e} l'arc AKC, a to\add{ute} la circonference de la terre, ou du cercle en supposant le seul mouuem\textsuperscript{t}. iournal\add{ier} ie dis que le graue ne descrira pas vn cercle, mais l'helice, ou la spirale AD, 21, \ill{} 15, 12. suyuant la figure marquée par points, dont voicy la proprieté, (supposant la proportion de la descente de Galilée) c'est que quelconque ligne droite EK men\add{ée} du centre, et coupante la spirale au point 21, et la circonference du cercle a\add{u} point K, et les droites EA, et EC, aussy menées, la droite EK sera à la droite K\ill{} en raison doublée de la circonference AKC, à la circonference AK. vne autre excellente proprieté est que l'espace de la spirale compris par AD, 21, E, \ill{} par la droite EA est au secteur AKCE comme 8 à 15. La premiere proprieté se prouue par la nature de la spirale, car l'arc AKC estant diuisé en autant de parties esgales qu'on voudra, et posé que la ligne BD parcourue par le graue d\ill{} la premiere partye de temps soit vne mesure, la ligne F26 parcourue en deux pr\ill{} temps sera de 4 mesures ; G25 de 9 ; H24 de 16 ; I23 de 25 ; K21 de 36, et ainsy \ill{}
autres suyuant
\note{Les \ill{} de ce paragraphe sont des fins de ligne cachées dans la reliure, non des mots mal écrits. « autres suyuant » est la réclame du bas de la page ; la phrase reprend en tête de la page suivante.}

\page{12}
\note{En haut à droite, « 6 ». La page continue sans interruption la remarque sur la p. 94 commencée au bas de la vue 11. Le tiers inférieur de la page est blanc ; l'écriture qu'on y voit en transparence est celle du verso.}

\note{Table dans la marge gauche, en regard du texte, en trois colonnes, les deux premières reliées par des points de conduite ; au-dessus de la troisième colonne, « … 65536 ». Sous la table, séparés par des traits, les totaux et la comparaison des produits croisés. Elle est donnée ici telle qu'elle est écrite.}
\[
\begin{array}{rr|r}
 & & 65536 \\
1 & 255 & 65025 \\
4 & 252 & 63504 \\
9 & 247 & 61009 \\
16 & 240 & 57600 \\
25 & 231 & 53361 \\
36 & 220 & 48400 \\
49 & 207 & 42849 \\
64 & 192 & 36864 \\
81 & 175 & 30625 \\
100 & 156 & 24336 \\
121 & 135 & 18225 \\
144 & 112 & 12544 \\
169 & 87 & 7569 \\
196 & 60 & 3600 \\
225 & 31 & 961 \\
256 & 0 &
\end{array}
\]

\marginal{nombres des quarrés de la derniere Colomne 592008 \quad nombre des mesmes quarrés sans le premier qui est 65536 : 526472 \quad l\add{e}d. quarré 65536 pris \struck{seize} fois 1048576 \quad plus grande : 592008/1048576 $\parallel$ 8/15 \quad moindre : 526472/1048576 $\parallel$ 8/15 \quad au 1\textsuperscript{er} plus grande : 8880120 $\parallel$ 8388608 \quad au 2\textsuperscript{me} moindre : 7897080 $\parallel$ 8388608}
\note{Le début des mots de gauche passe dans la reliure (« \add{le}d. quarré », « \add{pr}is »). Sous « seize » biffé est écrit un « 8 » ou un « 16 » mal formé, non tranché. « plus grande » et « moindre » sont écrits au-dessus des deux comparaisons, qu'ils qualifient. Le double trait vertical « $\parallel$ » oppose ici les deux termes d'une comparaison : les deux dernières lignes en sont les produits croisés ($592008 \times 15$ et $1048576 \times 8$, puis $526472 \times 15$ et $1048576 \times 8$). Les nombres de la table et des totaux ont été vérifiés : ils sont justes.}

autres suyuant l'ordre naturel des nombres quarrés, d'ou s'ensuit la premiere proprieté. La demonstration de la 2\textsuperscript{me} proprieté se fait en inscriuant, et circonscriuant a la spirale des secteurs de cercle, come sont les inscrits E3,4, EMN, E67, EPQ, E9 10, EST, E12 13, et cæt. Et les circonscripts E23, ELM, E56, EOP, E89, ERS, E11 12, E14 15, et cæt. ce qui se peut faire de sorte que la difference des inscripts aux circonscripts soit moindre que nulle autre quantité proposée, come a fait Archimede : voicy l'ordre des inscripts. Qu'il y ayt autant de nombres quarrés suyuant l'ordre naturel en commencant par 1, 4, 9 et cæt. comme il y a de l'arc \add{AKC} au secteur ou de secteurs AEB, BEF, FEG et cæt. c'est à dire 16 en nostre figure, et partant le nombre quarré sera 256 ; duquel ostés continuellem\textsuperscript{t}. les moindres a scauoir 225, 196, 169, 144 et cæt. ces nombres 31, 60, 87, 112, 135, 156 et cæt iusques à 255 demeureront ; les quarrés de tous ces nombres estans pris ensemble a scauoir 961, 3600, 7569 et cæt auec le quarré du nombre 256, a scauoir 65536 ; ont plus grande raison au mesme quarré 65536 pris 16 fois (suyuant le nombre des partyes esgales de l'arc AKC) que 8 à 15 : Et les mesmes quarrés 961, 3600 et cæt. pris ensemble sans le \struck{mesme} quarré 65536 ont moindre raison a ce mesme quarré 65536 pris 16 fois que 8 à 15 : ce qui est veritable en toute sorte de nombres de quarrés moindres ou plus grands que 16 : maintenant la droite CE estant posée de 256 mesures, la ligne 27,3 sera de 225 ; 28 M de 196 ; 29.6 de 169 et cæt. partant E3 costé du premier secteur tant inscript que circonscript sera 31, le costé du 2\textsuperscript{me} EM sera 60. celuy du 3\textsuperscript{me} E6 de 87 et cæt suyuant l'ordre des nombres precedent. Donc le premier secteur est au 2\textsuperscript{me} comme le quarré de 961 à celuy de 3600, parce que les secteurs sont entierem\textsuperscript{t}. semblables, et le 2\textsuperscript{me} sera au 3\textsuperscript{me} comme le quarré de 3600, à celuy de 7569 et cæt. Et la som\textsuperscript{e} de tous les secteurs circonscripts sera au secteur AEB pris 16 fois, c'est à dire à tout le secteur AKCE, en plus grande raison que 8 à \struck{\ill{}} 15 ; au lieu que la som\textsuperscript{e} des inscripts a moindre raison. Or les inscripts sont moindres que les circonscripts d'vne quantité moindre que telle petite quantité qu'on voudra, d'ou il s'ensuit que l'espace compris par la spirale, et la droite EA est au secteur AKCE, comme 8 à 15. toute la dificulté consiste à prouuer le lemme susd., touchant la som\textsuperscript{e} des nombres quarrés (comparée au nombre quarré 65536 pris 16 fois) produite par les nombres 31, 60, 87 et cæt. ; or l'on trouuera que ce lemme est veritable en en cherchant la demonstration.
\note{« AKC » est écrit au-dessus de la ligne, avec un signe d'insertion après « l'arc ». « le quarré de 961 à celuy de 3600 » est tel quel : ce sont déjà les carrés de 31 et de 60. « som\textsuperscript{e} » : « somme ». Un chiffre biffé devant « 15 ».}

Et bien que le secteur AKCE fut plus grand que le demicercle, ou mesme qu'il fut le cercle entier, et que l'helice eut fait son tour entier la demonstration seroit vraye, pourueu que l'on comparast l'espace de ceste spirale auec tout son secteur, ou auec tout le cercle, et mesme l'on peut considerer la mesme chose en plusieurs reuolutions d'helice. Et si le mouuem\textsuperscript{t}. ne commence pas dans l'equinoctial ou l'equateur, il se descrira vne spirale autour du cone dont l'espace sera a la surface du cone renfermée par ses bornes cor\textsuperscript{e} 8 à 15, ce qu'on demonstrera de la mesme facon.

p. 96 au droit de la ligne 12.
voyes l'helice descritte par m\textsuperscript{r}. fermat touchant lad. descente, laquelle est cy deuant demonstrée. sur la page 94 et 95.

lig. 19. apres ces mots (Il deuroit tumber en 6 heures)
de la surface de la terre iusques au centre

lig. 36.
Il a du depuis confessé par lettre qu'il n'auoit parlé de ce demicercle que par caprice et galanterie, et non a bon escient, et qu'il auoit preueu l'helice.
\note{« Il » : Galilée, d'après la remarque sur la p. 93 (vue 11).}

\page{13}
\note{Verso du feuillet 6. Aucun nombre dans les coins. Contre le bord droit de la vue se voit, en partie, la marge gauche de la page suivante (une table de nombres et « nota le… ») : elle est transcrite à la vue 14.}

p. 96. sur la fin
voicy vne autre demonstration, Que ceste ligne ne peut estre vne Circonference

\note{Figure en haut à droite : un quart de cercle de centre A, le rayon AC horizontal (E sur AC), l'arc passant par D et H ; un second cercle, plus petit, passe par A, I, K et C ; les droites AID, AKH, EK et KC, et une droite en pointillé prolongeant IK au-delà de K vers la droite ; le trait IK est repassé à l'encre forte.}

Que le point A soit le centre de la terre AC son demydiametre et que le cercle CHD represente l'equateur, les espaces que parcourt le graue en droite ligne, sont touciours en raison double de ceux qu'il fait dans le cercle, comme si le graue estant en C vient premierem\textsuperscript{t} en K, et puis en I menant la droite AKH, AID, la raison de HK à ID sera doublée de la raison de l'arc CH à l'arc CD, et si par exemple l'arc CD est double de l'arc CH, la droite ID sera quadruple de la droite HK. voyons si en supposant que le cercle descrit autour du diametre AC passe par les points IK la droite ID est quadruple de la droite HK lors que l'arc CD est double de l'arc CH. Menés la droite EK qui passe par le milieu de la droite AC, c'est a dire par le centre du cercle IKC, en sorte que KE soit esgale à AK, et qu'elle rencontre AC prolongée si besoin est en L. A cause de l'esgalité des lignes AK, KL, l'angle KLC sera esgal à l'angle KAL, ou son esgal IAH, mais les angles AIK, LKC sont aussy esgaux, et partant les triangles AIK, LCK sont semblables et d'autant que le costé KC est esgal au costé IK, le costé CL sera esgal au costé AI. D'ailleurs les triangles isosceles AKL, AEK ayants les angles sur les bases esgaux, ils seront semblables et le rectangle \struck{et le rectangle} LAE sera esgal au quarré de AK. soit AC, b : AK, e : CL, a : KH sera $b - e$ ; AL, $a + b$. le rectangle LAE, $\frac{1}{2}ab + \frac{1}{2}b^{2}$, lequel est esgal à ee qui est le quarré de AK et par consequent la droite CL ou son esgale AI sera $\dfrac{+ee \mp \frac{1}{2}bb}{\frac{1}{2}b}$ D'ou il s'ensuit que ID est $\dfrac{bb - ee}{\frac{1}{2}b}$, ce qui debueroit estre esgal à quatre fois $b - e$, qui est la valeur de KH, or si cela estoit $2be$ seroit esgal à $bb + ee$, ce qui n'est pas, mais au contraire $2be$ est moindre que $bb + ee$, d'autant que $b$ est plus grand que $e$, et que les quarrés de deux droites surpassent \struck{le} \add{leur} rectangle pris deux fois, du quarré de leur difference et par consequent ID est moindre que quatre fois KH. Donc la ligne que descrit le graue en descendant au centre n'est pas vne circonference de cercle. Or si ID est touciours moindre que le quadruple de KH, comme il est prouué, elle est aussy touciours plus grande que le double de \struck{KH} KH, car $\dfrac{+bb - ee}{\frac{1}{2}b}$ est plus grand que $2b - 2e$, à cause que $b$ qui est la marque de AC est plus grand que $e$ qui est la marque de AK. Mais on peut touciours faire que ID ayt à KH la proportion que l'on voudra pourueu qu'elle soit moindre que la quadruple et plus grande que la double. supposons que la proportion de $+s$ à $\frac{1}{2}b$ ayt ceste condition, que dessus, suyuant qu'on a trouué, ID, $\dfrac{bb - ee}{\frac{1}{2}b}$ et KH $b - e$ et par consequent ces quatre lignes seront proportionnelles
\[ \frac{bb - ee}{\frac{1}{2}b} \,/\, b - e \parallel s \,/\, \tfrac{1}{2}b \]
et le rectangle $bb - ee$ sera esgal au rectangle $+bs - ee$. D'ou ie tire ceste equation $-ee + es \parallel +bb$ \marginal{à $-bb + bs$ en transposant \struck{\ill{}} $bb - ee$ sera esgal à $+bb - ee$ \uncertain{ie ay}}, et d'autant que $s$ est plus grand que $b$, le 3\textsuperscript{me} terme sera vne quantité negatiue, co\textsuperscript{e} le premier et partant p\textsuperscript{r} scauoir la valeur \struck{\ill{}} de $e$ qui est AK, supposant dans la figure que AB soit $s$, FB $s - b$. AE, $b$. AK sera $e$. Apres cela inscriuant dans le cercle IKC la droite AK egale a AZ vous la prolongerés iusques en H, et faysant l'arc IK esgal a l'arc AC, vous tirerés la droite AID, cela fait, ie dis que ID est à KH co\textsuperscript{e} $s$ a $\frac{1}{2}b$ ; car puisque AZ, ou AK ou $e$ est de telle grandeur que $-ee + es$, \add{est} esgal \struck{$+bb - ee$} et ces quatre droites seront proportionnelles $\dfrac{bb - ee}{\frac{1}{2}b} \,/\, b - e \parallel s \,/\, \frac{1}{2}b$ or scauons que $+bb - ee$ est ID ; $b - e$, KH donc ID est à KH, co\textsuperscript{e} $s$ a $\frac{1}{2}b$ ce qu'il falloit demonstrer

Ceste ligne sera droite soubs le pole, helice plane sous l'equateur, et par tout ailleurs helice solide descrite sur la surface d'vn cone isoscele, qui a p\textsuperscript{r} base la parallele ou commence le mouuem\textsuperscript{t}, et p\textsuperscript{r} sommet le centre de la terre

\note{Notation : lettres minuscules pour les longueurs, $bb$, $ee$ pour les carrés (une fois $b^{2}$), « $\parallel$ », double trait vertical, pour « est à … comme » dans la proportion et pour l'égalité dans l'équation. Dans l'expression de AI, le signe entre $+ee$ et $\frac{1}{2}bb$ est un « $+$ » et un « $-$ » superposés, rendu $\mp$ ; la marge gauche récrit la même fraction « $+ee - \frac{1}{2}bb$ » sur « $\frac{1}{2}b$ ». Dans la première fraction de ID, un petit « b » est écrit au-dessus du signe moins. Dans « $-ee + es \parallel +bb$ », le « $bb$ » porte au-dessous un « $bs$ » plus petit, et un renvoi dans la marge gauche (un petit « x » surmonté d'un « a ») appelle la note marginale donnée ici comme note marginale, dont la fin est mal lue. Dans « \add{est} esgal \struck{$+bb - ee$} », un signe d'insertion est placé sous « esgal » sans que rien soit écrit au-dessus. « le 3\textsuperscript{me} terme » : le « 3 » est surchargé. Dans la marge gauche aussi, en regard du début de la démonstration : « AC $\parallel$ b », « AK $\parallel$ e », « CL $\parallel$ a », « KH $\parallel$ b $-$ e », « AL $\parallel$ a $+$ b », et un mot biffé.}

\note{Au pied de la page, à droite, la réclame « p. 98 au commencem\textsuperscript{t} ».}

\page{14}
\note{En haut à droite, « 7 ».}

\note{Dans la marge gauche, en haut, une petite table sans titre, deux colonnes de nombres séparées par un trait : 3 | 1 ; 15 | 8 ; 14 | 9 ; 45 | 32 ; 33 | 25 ; 91 | 72 ; 60 | 49 ; 153 | 128. Plus bas dans la même marge, une échelle verticale AG divisée en petits intervalles égaux par des traits, les points B, C, D, E, F marqués à droite de traits plus longs, et les nombres impairs 1, 3, 5, 7, 9, 11 à gauche, en regard du texte de la p. 100 ; les intervalles AB, BC, CD… croissent comme ces nombres. Ni la table ni l'échelle ne sont appelées dans le texte.}

p. 98 au commencem\textsuperscript{t}
Soit l'helice AFB descritte au dedans du cercle BCX en sorte qu'il y ayt touss\textsuperscript{rs}. mesme proportion de la circonference BCX à l'arc BC, que de AB à FC, ou que du quarré de AB au quarré de FC, ou que du cube de AB au cube de FC, ou bien de telle puissance de AB que l'on voud à la puissance du mesme degré de FC, donner la regle generale p\textsuperscript{r} trouuer la proportion du cercle BCX, à l'espace compris de la droite AB, et de ces helices AFB de quelque degré qu'elle puisse estre. la premiere helice est d'Archimede, et la seconde celle de Galilée p\textsuperscript{r} les graues, les autres ne tombent point en la pratique de la phisique, que l'on scache
\note{« touss\textsuperscript{rs}. » : « tousiours ». « voud » est tel quel. Figure à droite : un demi-cercle de diamètre XB, le point A sur le diamètre, une spirale partant de A, montant jusqu'en F et redescendant en B, et la droite AFC menée de A jusqu'au cercle en C.}

Au droit du corollaire
voyes le second point de la 1\textsuperscript{re}. obseruation

p. 99. \struck{ligne} au droit de la ligne 19.
voyes le 3\textsuperscript{me} point de la 1\textsuperscript{re} obseruation

p. 100.
p\textsuperscript{r} conceuoir le fondem\textsuperscript{t} de ce calcul et de ces tables, que le mobile A, soit vne minute a descendre de A en B, de sorte qu'estant en B il ayt acquis vne telle vistesse, que la continuant sans nouuelle augmentation il parcoure dans vn temps esgal à scauoir dans vne minute vn espace double de AB, si l'on adiouste encore pendant ce mesme temps vn espace esgal au premier produit par l'action de sa pesanteur, qui continüe touciours, l'on aura vn espace triple scauoir BC pendant le second temps esgal, c'est à dire pendant l'autre minute. Maintenant qu'il est en C et qu'il doibt faire vn espace double d'AC, pendant vn temps esgal à celuy qu'il est descendu en C, puis qu'il y est descendu en en deux minutes, il doibt en deux minutes faire deux fois l'espace AC, ou 8 fois l'espace AB par consequent durant vne minute quatre fois l'espace AB, parce que nous supposons que la vistesse acquise en C demeure touss\textsuperscript{rs} esgale, et sans croistre, et si l'on adiouste l'accroissem\textsuperscript{t} procedant de l'action de la pesanteur pendant ceste minute, qui adioustera vn espace esgal à AB, nous aurons en tout cinq espaces esgaux à AB p\textsuperscript{r} le troisiesme temps : semblablem\textsuperscript{t} puis que le mobile estant paruenu en D à acquis vne vistesse capable de luy faire faire vn espace double de AD, dans vn temps esgal à celuy qu'il est descendu en D, et puis qu'il est descendu en D en trois minutes, il fera en trois minutes deux fois l'espace AD, ou 9 fois AB, et pendant vne minute 6 fois \struck{AB} l'espace AB, auquel si l'on adiouste encores vn espace esgal à AB procedant de l'action de la pesanteur, qui continüe, l'on aura 7 espaces esgaux à AB p\textsuperscript{r} la quatriesme minute, et l'on aura la mesme chose des autres et cæt. ce qui est plus aisé à expliquer par les triangles, co\textsuperscript{e} est celuy de la 2\textsuperscript{me} proposition
\note{« en en deux minutes » : le mot est répété d'une ligne à l'autre. « ou 9 fois AB » : le « 9 » est repassé. Le paragraphe explique l'échelle de la marge.}

p. 101 au commencem\textsuperscript{t}
si les graues descendent par le mouuem\textsuperscript{t} de la matiere subtile autour de la terre, le fondem\textsuperscript{t} precedent ne sert pas. voyes les lettres de m\textsuperscript{r}. descartes sur ce subiet.

\marginal{nota lettres}

p. 102. au droit de la ligne 33.
parce que l'experience monstre qu'vne balle d'harquebuse employe du moins vne seconde minute apres les cent toises de sa sortie de blanc en blanc, il faudroit changer tout ce calcul p\textsuperscript{r} parler selon l'experience et p\textsuperscript{r} lors les nombres seront plus petits

p. 112 au droit du 1\textsuperscript{er} corollaire
Il dit dans ses dialogues de motu, qu'il a fait l'essay sur vn plan incliné long de 12 brasses, auec vne boule de metal

p. 120 au droit de la ligne 13.
Ce n'est pas vne demye ellypse excentrique mais elle en aproche seulement, c'est la trochoide, ou cycloide, qui enferme vn espace triple du \struck{cube} cercle qui la descrit./
\note{« Il » : Galilée. « nota lettres » est écrit dans la marge gauche, en regard de la remarque sur la p. 101.}

\page{15}
\note{Verso du feuillet 7. Aucun nombre dans les coins. Le bord droit de la page entre dans la reliure et cache la fin de plusieurs lignes ; les compléments sûrs sont entre crochets. À gauche, en transparence, la figure et l'écriture du recto.}

p. 121. au commencem\textsuperscript{t}
il faut corriger cecy suyuant l'onziesme obseruation mise à la fin
et au droit de la x\textsuperscript{me} proposition
Toute ceste proposition est bien mieux et plus ample dans vn traitté particulier à la fin du 3\textsuperscript{me} liure qui suit, c'est pourquoy il la faut oster d'icy, co\textsuperscript{e} inutile.

p. 125 au droit de l'onziesme proposition
Il faut encores voir sur cecy, la sixte obseruation
\note{« sixte » est lu sur un mot abrégé (« s\textsuperscript{te} »).}

p. 126. au droit de la derniere figure a la fin de la page
Apres auoir examiné ce que dit Salinas, l'on a trouué que ce qu'il dit de ceste diuision est faux, voyés le 1\textsuperscript{er} corollaire de la 18 prop. du liure de nos balistiques ou ceste ligne est coupée 12 fois en moyenne et extreme raison, et Salinas est entierem\textsuperscript{t} refuté
\note{« entierem\textsuperscript{t} » : le mot entre dans la reliure après « en ».}

p. 129. ligne 18 au droit (de ce qui est dans la parentese (qui pese 360 fois moins
Il faut corriger ce nombre en refaisant l'essay, auec vn trebuchet fort iuste, car la moelle de sureau ne pese que 193 fois moins que le plomb.
\note{La parenthèse ouverte avant « de ce qui » n'est pas fermée ; la seconde ouvre la citation du livre. Un signe d'insertion sous « nombre ».}

p. 130. au droit de la ligne 30.
voyes le 13\textsuperscript{me} et le 14\textsuperscript{me} article du 1\textsuperscript{er} liure des nouuelles pensées de Galilée, lequel tient dans sa response à Roca, que faisant descendre deux boules dans l'eau, l'vne de pierre et l'autre de plomb, de sorte que la pierre soit quatre fois plus pesante que l'eau, et le plomb trois fois plus pesant que la pierre, et par consequent 12 fois plus pesant que l'eau si l'eau est profonde de 1000 brasses, et que dans le vuide ces 1000 brasses se doibuent faire en 24\textsuperscript{$\prime\prime$}, l'eau empeschera le quart de la cheute, comme de la pesanteur, de la pierre et la douziesme partye de la cheute et pesanteur du plomb, de sorte que la pierre fera mille brasses d'eau en 30\textsuperscript{$\prime\prime$} et le plom en 26\textsuperscript{$\prime\prime$}, et si le plomb trouue son point d'egalité a cent brasses, la pierre à 90, desormais ils garderont touciours ceste mesme proportion ; de sorte que le plom ayant fait mille toises, la pierre en aura fait 900. De mesme dans l'air si la pierre est mille fois plus pesante que l'air, et le plomb 3000 fois plus pesant que le mesme air, il diminuera seulem\textsuperscript{t} la pesanteur de la vistesse de la pierre d'vne millesme partye, et celle du plomb d'vne troismilliesme, de sorte que le plomb tombant de cent brasses ne doibt retarder le mouuem\textsuperscript{t}. qu'il auoit dans le \struck{\ill{}} vuide, c'est au 3\textsuperscript{me} en raison doublée des temps, que de $\frac{1}{3000}$ partye de cent brasses, qui n'est qu'en vn mouuem\textsuperscript{t} vn pouce. et si vne boule de liege est cent fois plus legere que le plomb, parce que 100 est 30 fois en 3000, l'air ostera vne 30\textsuperscript{me} partye de la pesanteur, et de la vistesse du liege. voyés la 1\textsuperscript{re}. obseruation a la fin de ce volume, ou tout ce discours de Galilée a grand besoing d'examen, ce qu'on peut f\textsuperscript{re} par la raison et l'experience, voyez l'experience contre Galilée a la fin de l'article 15 de son premier liure.
\note{« quatre fois plus pesante que l'eau », « l'vne de pierre », « plus pesant que l'eau », « se doibuent faire » : fins de ligne complétées au bord de la reliure. Le passage « c'est au 3\textsuperscript{me} … qu'en vn mouuem\textsuperscript{t} vn pouce » est lu tel quel ; son sens n'est pas clair, et « au 3\textsuperscript{me} » peut couvrir un mot perdu dans la reliure. « f\textsuperscript{re} » : « faire ».}

p. 131. au commencement
Il est probable que les graues ne continuent pas a augmenter leur vistesse iusques au centre de la terre, et qu'apres quelque espace ilz ne l'augmentent plus, ou qu'ilz l'augmentent beaucoup moins, peutestre suyuant les approchem\textsuperscript{t}. des assymptotes
\note{« iusques », « qu'ilz » : fins de ligne à demi cachées dans la reliure.}

p. 135 au droit de la 15\textsuperscript{me} proposition
Quelqu'vn scauoir m\textsuperscript{r}. decousu dit auoir experimenté que la chorde ne doibt auoir que deux pieds et 10 pouces p\textsuperscript{r} faire 3600 tours dans vne heure, il est libre a chascun d'en adiuster la longueur trois pieds suffiront
\note{« auoir » : fin de ligne à demi cachée.}

p. 136. au commencem\textsuperscript{t}.
1\textsuperscript{o}. vn baston rond ou quarré pendu par le bout doibt estre plus long que le filet isochrone, et ce en raison sesquialtere proxime : ou suyuant le P. fabry comme l'aire du demicercle au quarré de son rayon : et dit qu'vn graue tombera seulem\textsuperscript{t} aussy viste que led. baston fera sa demie vibration.

2\textsuperscript{o}. les triangles isosceles presque depuis 20 iusques à 60 degrés, pendus par la pointe sont au filet isochrone, presque comme 4 à 3 et par la base co\textsuperscript{e} 3 à 2 neantmoins celuy de 60 degrés est co\textsuperscript{e} 6 à 5 par la pointe

3\textsuperscript{o}. le Triangle rectangle tant par la pointe que par la base est esgal au filet
\note{Dans la marge gauche, en regard des trois articles, les chiffres « 1\textsuperscript{o} », « 2\textsuperscript{o} », « 3\textsuperscript{o} », qui répètent ceux du texte. Au pied de la page, à droite, la réclame « 4\textsuperscript{o} le Triangl\add{e} ».}

\page{16}
\note{En haut à droite, « 8 ». La page continue la liste commencée au bas de la vue 15 ; dans la marge gauche, les chiffres « 4\textsuperscript{o} », « 5\textsuperscript{o} », « 6\textsuperscript{o} », « 7\textsuperscript{o} » répètent ceux du texte. En haut à droite, une esquisse très pâle, sans lettres : un triangle isocèle suspendu par la pointe à un fil vertical terminé par une flèche.}

4\textsuperscript{o}. le triangle de 110 degrés par la pointe co\textsuperscript{e} 2 à 3. par la base co\textsuperscript{e} 8 à 15
celuy de 120 degrés par la pointe et la base quasi de mesmes
celuy de 135 \add{degrés} par la pointe comme 2 à 5 par la base comme 1 à 4
celuy de 150 degrés par la pointe co\textsuperscript{e} 1 à 4 et par la base co\textsuperscript{e} 1 à 12
celuy de 170 degrés par la pointe co\textsuperscript{e} 1 à 31 proxime et par la base co\textsuperscript{e}. 1 à 36
\note{« degrés » après « 135 » est écrit au-dessus de la ligne.}

5\textsuperscript{o} Cercle, son diametre est au filet co\textsuperscript{e} 7 à 5 proxime
Demicercle tant par le conuexe que par la base comme 5 a 6
Circonference egal. Demye circonference par le conuexe egal au rayon par la base le filet est sesquialtere
Quart de cercle par la pointe co\textsuperscript{e} 3 à 2 par la base co\textsuperscript{e} $1\frac{5}{7}$
\note{Le dernier rapport, « $1\frac{5}{7}$ », est écrit sans second terme.}

6\textsuperscript{o} le Cone de 20 degrés par la pointe co\textsuperscript{e} 6 à 5 par la base co\textsuperscript{e} $2\frac{1}{4}$ à I.
La sphere co\textsuperscript{e} $1\frac{4}{9}$ à vn presque co\textsuperscript{e} le Cercle
La demye sphere co\textsuperscript{e} 11 a 16 par la pointe et la base
La parabole dont la base est sesquialtere de son axe, par la pointe co\textsuperscript{e} 9 à 7 et par la base co\textsuperscript{e} 18 à 11.
\note{« 16 » est souligné.}

7\textsuperscript{o} la roulette, ou trochoide par la pointe et la base co\textsuperscript{e} 3 a 2
parabole 2\textsuperscript{me} de Toricelli p\textsuperscript{r} boyre esgalem\textsuperscript{t}, esgale au filet par la pointe et la base.

p. 137 proche la 3\textsuperscript{me} ligne
voyés depuis la 23\textsuperscript{me}. proposition du liure de la voix, iusques à la 29\textsuperscript{me}, p\textsuperscript{r} les proprietés de l'ellypse.

p. 142. a la ligne 39 en ces mots (co\textsuperscript{e} 360 à 1,
Ayant fait l'essay plus exact auec vn trebuchet qui trebuche auec $\frac{1}{16}$ partye de grain, i'ay trouué que la boule de moelle de sureau est à celle de plomb comme 1 à $128\frac{2}{3}$.

p. 156. au droit de la 6\textsuperscript{me} ligne.
Il est douteux si cela arriue aux bales, il semble que no, et que la fabrique des arbalestes et des fleches soit cause que l'on ne tire pas horizontalem\textsuperscript{t}.

\marginal{au 3\textsuperscript{me} liure du mouuement}

p. 157. prop. premiere
voyés le 10\textsuperscript{me} Apiarium progymn. 2\textsuperscript{o} prop. 3\textsuperscript{a} ou par l'ellypse il est monstré de combien les triangles, et partyes de la chorde sont plus grandes lors qu'on touche les chordes plus pres du cheualet, et qu'on luy fait faire autant de chemin, que lors qu'on la touche au milieu entre le cheualet et le sillet, ou en tel autre lieu qu'on voudra. Il est certain que sy l'on tire la chorde AB en telle facon qu'elle vienne de D en E et en F et cæt. p\textsuperscript{r} aller en C, elle sera touciours de mesme longueur, supposé que AFEB est vne ellypse. mais il faudroit supposer que les focus de l'ellypse sont A et B, et qu'on cognust la resistence de tous les points et tensions de la chorde.
Si l'on met vne ligne parallele à la chorde de luth ascauoir CD pour lors la chorde BEA sera la plus alongée car le triangle Isopleure BFA, est le plus court en peripherie de tous ceux qui sont entre mesmes paralleles quoyque le triangle BFA soit esgal en aire du triangle BEA, et à tous les autres possibles, qui sont entre mesmes paralleles et ont touss\textsuperscript{rs} p\textsuperscript{r} base AB de sorte qu'icy il y a esgalité d'aires, et inegalité de circuits, et en l'ouale esgalité de circuit et inegalité d'aires, et tant entre paralleles, que dans l'ellypse, le triangle Isopleure est le plus petit in ambitu./
\note{« au 3\textsuperscript{me} liure du mouuement » est écrit dans la marge gauche, sous un trait qui la traverse, en regard de « p. 157 » : la remarque passe au livre suivant. Deux figures accompagnent la remarque. La première, marquée « figure 1\textsuperscript{re} », à droite : une demi-ellipse de grand axe AB, C et D sur l'axe, F et E sur la courbe, avec FC perpendiculaire à l'axe et les droites AE, EB, ED, FD, FB. La seconde, marquée « figure 2\textsuperscript{me} », à gauche : la droite CD parallèle à AB, E et F sur CD, les triangles AEB et AFB. Dans la figure 1\textsuperscript{re}, un trait relie le haut de la courbe au mot « voudra » du texte. « no » : « non ».}

p. 160 prop. 2\textsuperscript{me}
Sy l'on met icy la figure du dernier corollaire de la 18\textsuperscript{me} proposition du 3\textsuperscript{me} liure des instrumens, l'on n'aura pas besoing de changer le D en H./

\page{17}
\note{Verso du feuillet 8. Aucun nombre dans les coins. Le bord droit de la page entre dans la reliure ; les fins de ligne qu'on devine sont complétées entre crochets, et ce qui manque est laissé en lacune. Contre ce bord se voit la marge gauche de la page suivante (une échelle de syllabes, des figures), transcrite à la vue 18.}

\struck{\ill{}} En la mesme page 160. ligne antipenultiesme propos. 2\textsuperscript{me}
Il faudroit changer toute ceste proposition, afin de dire que la chorde AC bandée en D, reuient par la force de quelque esprit interne, ou par celle de la matiere subtile, et que ce mouuem\textsuperscript{t} accompagne touss\textsuperscript{rs}. la chorde iusques a D, et parce que la fleche qui seroit appuyée sur la chorde ADC, au point D, ne quitteroit point la chorde qu'en arriuant au point B, il s'ensuit que la chorde va du moins aussy viste en B co\textsuperscript{e} en D, or p\textsuperscript{r} trouuer la proportion de la vistesse de D en B, il faut supposer 1\textsuperscript{o} que le mouuem\textsuperscript{t}. imprimé \add{a} \struck{en} D pour retourner a B ne cesseroit point d'aller d'esgale vistesse, si elle se faisoit dans le vuyde, et parce que l'air resiste, et retarde la force interne de D renforce touciours et recompense le retardem\textsuperscript{t}. or il est certain que sy auec le mouuem\textsuperscript{t}. imprimé continuant d'egale vistesse, la force interne qui l'accompagne luy imprimoit encores vn degré de vistesse en chasque temps la vistesse de D allant en B croistroit en raison doublée des temps co\textsuperscript{e} 1, 3, 5, 7 et cæt sy le mouuem\textsuperscript{t}. se faisoit dans le vuyde, mais parceque l'air resiste, et que la force interne qui pousse, se diminüe vniformem\textsuperscript{t} qu'elle aproche de B, il s'ensuit que DB se mouuant de D à B, son mouuem\textsuperscript{t}. s'augmente touciours à raison de la force interne qui l'accompagne touss\textsuperscript{rs}, mais p\textsuperscript{r} entendre ceste vistesse soit tout le temps du retour diuisé en 88 partyes, si au premier temps l'arc ou la chorde ADC, c'est à dire D fait 8 parties ; au 2\textsuperscript{me} temps, 20 ; au 3\textsuperscript{me}, 28 ; et au 4\textsuperscript{me} 32, car au lieu de faire au deuxiesme temps deux fois 8 la vistesse se diminüant d'vn et ne fait plus que 15, mais parce qu'au deuxiesme temps l'augmentation de la vistesse doibt estre 5, D fera 20 espaces.
\note{Le premier mot de la remarque, avant « En », est biffé et illisible. « a » est écrit au-dessus de « en » biffé. Figure à gauche, en haut : un triangle ADC (A et C aux extrémités de la base, D au sommet), la hauteur DB prolongée au-dessous de la base par une flèche ; la hauteur est divisée en quatre, numérotée 1 à 4 de B vers D à gauche, et 1 à 4 de D vers B à droite. Plus bas, un losange en pointillé ACBD de diagonale AB (E au milieu), C en haut, D en bas. Les deux « au 2\textsuperscript{me} temps, 20 » et « D fera 20 espaces » s'accordent ; la somme 8, 20, 28, 32 fait bien 88.}

p. 165, a la proposition 4\textsuperscript{me}
voyés la raison de ce retour, et des ressorts à la 2\textsuperscript{me} page de la preface generale, qui consiste en ce que le torrent de la matiere subtile, qui a ses partyes rondes, ne trouuant plus les pores des corps plus ronds mais en ouale elle \struck{pousse} presse les bords estroits des cors ou pores p\textsuperscript{r} passer plus aysem\textsuperscript{t} a trauers les corps, en redressant leurs pores en rond. l'on peut aussy dire qu'il y a quelques esprits internes dans la chorde tendue qui la font retourner a son assiette.
\note{« la 2\textsuperscript{me} page de la preface generale » : renvoi au livre. « pousse » est biffé et « presse » écrit à la suite. « l'on peut », « retourner » : fins de ligne à demi cachées. Figures à droite du texte : un demi-cercle dont le diamètre porte D, avec C au-dedans et deux cordes ; un triangle de sommet F et de base E ; plus bas, une figure en forme de lentille, de diamètres CD (vertical) et horizontal se coupant en E.}

\struck{prop. 166 p.} pag. 166. prop. 5 il y a ces deux figures

p. 168. a l'aduertissem\textsuperscript{t}.
ce traitté ne s'estant pas trouué prest, a esté remis et differé apres la 24\textsuperscript{me}. propos. de ce liure c'est à dire iusques à la fin

pa. 169. lig. 22
voyés la 4. prop. du 4\textsuperscript{me}. article du 1\textsuperscript{er} liure latin de l'harmonie, iointe à l'hydraulique et l'experience faite depuis pour monstrer que la corde bat l'air 108 fois, quand elle est à l'vnisson auec vn tuyau d'orgue de 2 pieds bouché.
\note{« l'vnisson », « elle est » : lecture au bord de la reliure. « 108 » peut aussi se lire « 128 ».}

pa. 171. prop. 7\textsuperscript{me} il y a ceste figure

pa. 173 au commencem\textsuperscript{t}.
Galilée tient que la chorde bandée qui se laisse faire vne figure parabolique, dont elle doibt faire la mesme figure dans ses tours et retours
\note{Lecture littérale ; « se laisse faire » est incertain, et un mot peut manquer au bord de la reliure après « bandée qui ».}

lig. 22 apres ces mots (soubsdouble de BA)
par supposition car dans la figure, AB, est plus \struck{grande} que double de RS

lig. 35 apres ces mots (DE est plus long que HV.)
l'on peut adiouster que DE donneroit plus d'impetuosité, bien que HV allast aussy viste, à raison du plus long chemin de DE, qui semble seruir à vne plus grande impression./

lig. 3 depuis la fin apres ce mot (proportionnées)
Toutes les experiences me font conclure que la force de l'arc ou du coup doibt estre 4 fois plus grande p\textsuperscript{r} enuoyer les missiles deux fois plus loing, et ainsi de suite en raison doublée des proiections./
\note{« estre » et « de suite » sont suppléés au bord de la reliure. Contre ce bord, en regard de « DE est plus long que HV », on lit de travers « ce nombre faux », qui est de la marge gauche de la page suivante.}

\page{18}
\note{En haut à droite, « 9 ».}

\note{Dans la marge gauche, en regard de la remarque sur la p. 185, une table des poids qui tendent la corde pour les huit degrés de l'octave, écrite en deux colonnes séparées par un trait :}
\[
\begin{array}{l|l}
\text{vt} & 24.\ \text{Liures} \\
\text{si} & 21\,\frac{3}{16} \\
\text{la} & 16\,\frac{2}{3} \\
\text{sol} & 13\,\frac{1}{2} \\
\text{fa} & 10\,\frac{2}{3} \\
\text{mi} & 9\,\frac{1}{8} \\
\text{re} & 7\,\frac{33}{81} \\
\text{vt} & 6.\ \text{liures}
\end{array}
\]
\note{Valeurs données telles qu'elles sont écrites. Sous la table, deux figures : un demi-cercle de diamètre EG, avec F au sommet et la perpendiculaire FH ; un demi-cercle de diamètre AC, avec B au sommet et la perpendiculaire BD — les cordes EFG et ABC de la remarque sur la p. 189. Plus bas, la figure du triangle KNP (voir cette remarque).}

p. 183. deuant le Corollaire immediatem\textsuperscript{t}.
Il faut neantmoins touciours auoir esgard aux vent qui se trouuent souuent dans les metaux, et a ceux qui sont plus ou moins battus, par exemple le plomb batu ne tient pas tant d'espace que le fondu

p. 185. au droit de la table qui gist
Supposé que le 1\textsuperscript{er} poids de six liures face bander la chorde p\textsuperscript{r} faire l'vt graue, les autres poids suyuants la tendant l'vn apres l'autre la feront monter aux autres 7 sons compris dans \struck{l'octaue}, l'octaue, car les raisons de ces poids sont doublées des raisons des interualles diatonics de sorte que 8 cordes de mesme longueur et grosseur tendües par ces poids feront faire les degrés de l'octaue de l'vt sol \uncertain{fa} vt par b.
\note{« l'octaue » est écrit, biffé, puis récrit. La fin de la dernière ligne, « de l'vt sol \uncertain{fa} vt par b », est lue lettre à lettre ; le sens n'en est pas sûr.}

p. 189. prop. XV\textsuperscript{me}
Pour entendre la raison de ceste tension quadruple, que les chordes doiuent auoir, supposons que les chordes ABC, et EFG soyent en tout esgales, sinon qu'ABC soit plus tendüe que EFG, de sorte qu'elle monte à l'octaue, et qu'elles soyent esgalem\textsuperscript{t}. esloignées de leurs lignes de direction, c'est a dire que BD, et FH soyent esgales, il est certain qu'il ne faut ny plus ny moins de force et de temps, en contant l'vn auec l'autre, p\textsuperscript{r} faire que ABC, reuienne iusques à D, que p\textsuperscript{r} f\textsuperscript{re} que EFG reuienne iusques \struck{H} a H, c'est à dire que sy ABC a plus de force, il luy faudra moins de temps à proportion, car toutes les autres choses estants esgales, ceste inegalité de la force ne peut estre recompensée que par celle du temps. Il est certain aussy que puisque ABC fait l'octaue sur EFG, elle n'employe que la moitie d'autant de temps à passer de B à D, que EFG à passer de F à H, si bien qu'il ne reste plus qu'a scauoir combien la force qui la meut doibt estre plus grande que celle qui meut l'autre, afin que ceste force et ce temps contés ensemble facent en toutes deux la mesme somme. or parce que la force agit touciours esgalem\textsuperscript{t}, et que l'impression qu'elle fait à chasque moment demeure iusques à la fin du mouuem\textsuperscript{t}, on peut representer le temps par vne ligne, co\textsuperscript{e} KL ou KN, et la force par vne autre, co\textsuperscript{e} NO, LM, ou NP, en sorte que l'vne et l'autre ensemble soit representé par le triangle KNO, ou KLM, ou KNP, à scauoir puisque ABC n'employe que la mortye d'autant de temps à aller de B à D, que fait EFG à aller de F à H, Je represente le temps de ABC par KL prise à discretion, et celuy de EFG, par KN, que ie fais double de KL, puis ie represente la force de EFG par NO prise de rechef à discretion, et celle d'ABC par NP en vn temps esgal, et par LM, en vn temps moindre de moitié, et ceste LM, doibt estre telle, que le triangle KLM, soit esgal au triangle KNO, mais a cest effect elle doibt estre double de NO, ensuitte NP doibt estre quadruple de NO, donc la force qui meut ABC, doibt aussy estre quadruple de celle qui meut EFG, car lors qu'elles sont considerées d'elles mesmes, sans esgard de aucun \struck{temps}, elles ont mesme raison l'vne à l'autre que lors qu'elles sont considerées au regard du temps esgal.//
\note{« mortye » : « moitié », tel qu'écrit. Le mot biffé après « aucun » est couvert d'une tache ; on devine « temps ». Figure à gauche : un triangle KNP (K au sommet, N et P aux extrémités de la base), LM parallèle à la base, O sur la base, et la droite KO.}

p. 193. prop. 16\textsuperscript{me}
voyés la xvj\textsuperscript{me}. prop. du 1\textsuperscript{er} liure des instruments, et la 19\textsuperscript{me} du 3\textsuperscript{me} et le 1\textsuperscript{er} et 2\textsuperscript{me} dialogue de Galilée ou il dit qu'vne chorde se romproit par son propre poids, si elle estoit si longue qu'elle pesast autant que le poids qui la rompt, voyés les experiences du 2\textsuperscript{me} article de ses nouuelles pensées. Il dit que la resistence des cylindres est triple de la raison des diametres de leurs bases, et partant qu'vn baston soubsdouble en grosseur est huit fois plus aysé à rompre. Donc puis que la chorde de letton, dont le diametre de la base a $\frac{1}{6}$ de ligne rompt par la force de 18 liures, le cylindre ou la colomne qui a \struck{\ill{}} vn pied en diametre rompra seulem\textsuperscript{t} auec 13436928
\marginal{ce nombre me semble faux}
\note{« 13436928 » est lu chiffre à chiffre à fort grossissement ; la phrase s'arrête sur ce nombre, sans unité. La note « ce nombre me semble faux » est dans la marge gauche, en regard, de la main et de l'encre du texte.}

p. 197. coroll. 2\textsuperscript{me} ligne 6. apres ces mots (qu'vn degré de vistesse au premier pied.
Il faut changer cecy, car a la fin du 4\textsuperscript{me} pied il n'a que deux degrés de vistesse, parce que s'il n'augmentoit plus sa vistesse, à la fin du 4\textsuperscript{me} pied en continuant auec la vistesse acquise il feroit deux fois autant de chemin, ascauoir 8 pieds en mesme temps qu'il en a fait quatre./

p. 211 coroll. premier
Le S\textsuperscript{r} decousu dit auoir trouué, qu'elle ne doibt auoir que deux pieds et 10 pouces chascun la peut adiuster au plus pres, et la longueur de trois pieds sera bonne/

\page{19}
\note{Verso du feuillet 9. Aucun nombre dans les coins. À gauche, en transparence, la figure du triangle KNP du recto ; à droite de la remarque sur la p. 74, un triangle pâle, sans lettres, qui est aussi la figure du recto vue par transparence.}

p. 212 lig. 26.
voyez l'vsage de l'horloge vniuersel

p. 215. lig. 23.
Si en parlant on est esloigné de 39 ou 40 toises du fossé, la muraille de la ville respondant monstrera que le fossé a 20 ou 21 toises

p. 216. lig. 4.
voyés la 9 prop. de l'vtilité
\struck{p.} lig. 17. \struck{voyés}
voyés la resolution de cecy dans la 9\textsuperscript{me} propos. de l'vtilité et dans la 3\textsuperscript{me} obseruation.

p. 217 lig. 7.
m\textsuperscript{r}. descartes a eu vn echo en son iardin qui respondoit co\textsuperscript{e} le coq d'vn poullet à toutes sortes de bruits et de voix, ce qui venoit peutestre des brins d'herbe de la hauteur d'vn homme, dont la rencontre redoubloit les batemens
\note{« coq d'vn poullet » est lu tel quel ; comparer la remarque sur la p. 50 (vue 10), où le même écho répond « comme le cri d'vn poullet ». « peutestre » est lu sur un mot contracté.}

coroll. 4\textsuperscript{me}
La voix fait 1380 pieds dans le temps d'vne seconde / p\textsuperscript{r} suppleer a ce qui est icy voyés la 1\textsuperscript{re} obseruation du liure de l'vtilité.

p. 220. lig. 16.
La solution \struck{de} c'est que le fondroit va touciours d'vne esgale vistesse, et que le reflex est plus lent, voyez la 9 propos. de l'vtilité et la 3\textsuperscript{me} obseruation./
\note{« fondroit » : « foudroit », sans doute (la foudre, le tonnerre) ; lu tel qu'il est tracé. « le reflex » : fin de ligne au bord de la reliure.}

p. 222 lig. 35.
Elle n'en employe que 12\textsuperscript{$\prime\prime$} à monter et autant à descendre, et partant elle ne monte que 288 toises, voyés la 8 et 9 propos. de l'vtilité.
D'autres experiences m'ont monstré que les flesches montent en moins de temps qu'elles ne descendent, et que celle qui monte en 3 secondes descend en 5 secondes, par ainsy il faut corriger ce qui est icy et dans la 2\textsuperscript{me}. obseruation phisique et ailleurs

\marginal{xa.}
p. 223. lig. 32.
Galilée ne demonstre pas cecy, en son 3\textsuperscript{me} dialogue du mouuem\textsuperscript{t}, mais on le trouue demonstré dans les mechaniques de m\textsuperscript{r}. \uncertain{gaulion}.
\note{Le nom est tracé « gaulion » ou « goulion ». La marque « xa. », dans la marge gauche, est répétée à la remarque suivante ; son sens n'est pas donné.}

\marginal{xa.}
p. 224. deuant le corollaire premier
Ce liure est maintenant imprimé et traduit en francois en 4 liures.

p. 225. coroll. 2\textsuperscript{me}
Il faut corriger tout cecy puisque les missiles employent plus de temps à descendre qu'à monter

p. 226 ligne 5.
cecy est obscur, il faut l'esclaircir et l'expliquer.

\marginal{Du 1\textsuperscript{er} liure de la voix}

p. 33. \struck{lig} au droit de la figure
voyes la 31\textsuperscript{me}. page de l'vtilité de l'harmonie a la marge
\note{La note de la marge gauche, sous un trait qui la traverse, marque le passage à un nouveau livre ; la pagination du livre repart ici de 33.}

p. 74, lig. 41. apres ces mots (des mechaniques)
ces 8 poids en progression triple sont. / 1. 3. 9. 27. 81. 243. 729. 2187 /

p. 80 lig. 15.
plus\textsuperscript{rs}. nient qu'il s'ensuiue qu'vne chose doiue auoir toute sorte de perfections infinie encore qu'elle soit de soy mesme et independante, no seulem\textsuperscript{t}. parceque la nature de chasque chose est d'vne eternelle verité, qui ne depend point d'ailleurs, selon quelques vns mais aussy parce que la conception de l'independant n'est pas celle de la perfection infinie.
\note{« no » : « non ». « n'est pas » : la fin de ligne entre dans la reliure.}

Lig. 38.
M\textsuperscript{r}. de villiers a obserué l'année 1638 qu'au crane d'vn sourd et muet, des trois osselets, qui se trouuent en l'oreille, à scauoir l'incus, le malleolus, et le stapes, le premier manquoit en chasque sinus de chasque oreille, ce qui est tres notable p\textsuperscript{r} scauoir la cause de la surdité. /
\note{« stapes » : lecture au bord de la reliure, où l'on voit « sta ». « 1638 » est lu à fort grossissement.}

\marginal{Du liure second des chants}

p. 95. au droit de la musique
chant ancien des P\textsuperscript{res}. feuillants, qui chantoient les pseaumes et les litanies sur le chant qui suit, dans lequel le 2\textsuperscript{me} couplet varie.

\note{Portée de cinq lignes, clé d'ut sur la quatrième ligne, notation blanche, avec des barres de mesure et deux points d'orgue ; sous les notes : « Dixit dominus domino meo, sede à dextris meis Donec ponam inimicos meos scabellum pedum tuorum ». Les dernières lettres entrent dans la reliure. La musique n'est pas transcrite.}

\page{20}
\note{En haut à droite, « 10 ». On est toujours dans le « liure second des chants » annoncé en marge à la page précédente.}

p. 107. prop. 8\textsuperscript{me}
voyés le probleme arithmetique à la fin du 2\textsuperscript{me} tome de Centro grauitatis de guldinus, qui conte le nombre des lignes qui sont en toutes les dictions qui se peuuent faire des 23 lettres de l'alphabet, et combien tout cela feroit de volumes, et de bibliotheques, encores qu'en chasque diction nulle lettre ne se repetast ce qui se raporte a nre\textsuperscript{e} prop. 11 qui suit : or il trouue autant de lettres. 154600749126720214790543\uncertain{3}, dont le 1\textsuperscript{er} chiffre 1, vaut vn septilion, et le nombre des dictions qui contiennent ce nombre de lettres est 702730673303300980911155, dont le 1\textsuperscript{er} chyphre 7, vaut septante ou sept diz\uncertain{aines} de sextilions
\note{« lignes » est lu sur un mot abrégé (« lig\textsuperscript{s} » ou « lres ») et peut valoir « lettres ». « nre\textsuperscript{e} » : « nostre ». Le dernier chiffre du premier nombre est suivi d'une virgule serrée et se lit mal ; les deux nombres sont lus chiffre à chiffre à fort grossissement. « diz\uncertain{aines} » : le mot est tracé « dipaines » ou « dizaines ».}

p. 110. au droit du nombre 64\textsuperscript{me} de la table
Il y a 90 caracthers dans le 64\textsuperscript{me} nombre, de sorte qu'il y a au premier ternaire 221 huitilions, il y a aussy, 90 caracthers dans le liure de la grenze page 53 qui monstre combien vn monde d'autant plus grand que cestuicy, que cestuicy iusques iusques au firmament est plus grand qu'vn grain de sable, contiendroit de grains de sable en toute sa solidité.
\note{« la grenze » est lu tel quel ; c'est peut-être un nom (« de la Grenze » ?) ou un titre mal compris par le copiste. « iusques iusques » : le mot est répété.}

Mais si l'on prend vn nombre quarré magique, et que l'on veuille scauoir en combien de manieres il se peut varier, demeurant touss\textsuperscript{rs} magique, le nombre en sera encores plus grand, particulierem\textsuperscript{t}. si le nombre est vn peu grand, comme celuy de 22, dont vne enceinte pouuant estre ostée en sorte que les rangs 20 qui restent soyent encores esgaux en leur somme la variation n'en peut estre exprimée que par le nombre qui vient de la multiplication de 1332756462138013445206441942056960
par 11884960861556899184640000000000
et du produit par 114115536202195131378853478400000
et le quarré de 14, duquel on oste chasque enceinte iusques à 4, ce qui est du 4\textsuperscript{me} restant touss\textsuperscript{rs}. de mesme se peut varier en 33710564646400.
Quant a la variation absolüe, ou les bons sont auec les mauuais, c'est à dire les magiques auec les non magiques il ne faut que prendre la combination ord\textsuperscript{re}. du quarré par exemple de 9 qui est le premier, de 14, de 22 et ainsy des autres. et si l'on exprimoit toutes les varietés magiques du quarré de 22, il faudroit vn nombre de plus de 400 caracthers
\note{Les trois grands nombres sont écrits en colonne, l'un sous l'autre, et lus chiffre à chiffre ; « et du produit par » introduit le troisième. Leur exactitude n'est pas vérifiée. « ord\textsuperscript{re}. » : « ordinaire ». « le premier » : le premier carré magique qu'on puisse, selon le texte, varier ainsi.}

p. 111. au commencem\textsuperscript{t}.
p\textsuperscript{r} faire toutes les varietés sans se confondre, il faut tenir vn ordre à arranger le nombre des choses proposées : par exemple s'il y en a 4, come DIEU. Il faut retenir la 1\textsuperscript{re} lettre et changer l'ordre des trois autres, ie change donc les trois IEU en 6. facons, et parce que I est le premier, ie le retiens encores, et change EU en toutes les sortes qui sont 2, scauoir EU, et UE, puis ie change E p\textsuperscript{r} auoir EI, et puis ie mets les 2 autres en leurs 2 facons : puis, apres D ie mets la 4\textsuperscript{me} lettre U, et change encores les 2 autres en leurs 2 facons, et parceque la 4\textsuperscript{me} est la derniere, et qu'on ne peut plus en mettre d'autre au 2\textsuperscript{nd} lieu, ie change la 1\textsuperscript{re} D p\textsuperscript{r} mettre la 2\textsuperscript{de} I en sa place, et en suitte les 3 autres selon leur ordre, scauoir la 1\textsuperscript{re} au 2\textsuperscript{nd}, et puis la 3\textsuperscript{me} et 4\textsuperscript{me}, et ainsy I demeurant au 1\textsuperscript{er} lieu on fait les 6 changem\textsuperscript{ts}. des trois autres DEU. come deuant ce qu'estant accompli, on mettra la 3\textsuperscript{me} lettre E au 1\textsuperscript{er} lieu p\textsuperscript{r} f\textsuperscript{re} encores les 6 variations des trois autres, et enfin on mettra la derniere U au premier lieu p\textsuperscript{r} arranger les 3 autres DIE en leur 6 facons.
S'il y auoit 5 lettres differentes, il faut f\textsuperscript{re} co\textsuperscript{e} deuant les 24 changem\textsuperscript{t} des 4 dernieres, la 1\textsuperscript{re} demeurant touss\textsuperscript{rs} en son 1\textsuperscript{er} lieu, et puis l'on oste la 1\textsuperscript{re}, et l'on met en son lieu la 2\textsuperscript{de}, p\textsuperscript{r} f\textsuperscript{re} encores 24 variations et ainsy de la 3\textsuperscript{me}, 4\textsuperscript{me}, et 5\textsuperscript{me} lettre.
et s'il y auoit 6 choses co\textsuperscript{e} icy, 6 notes, on doibt faire le changem\textsuperscript{t} des 5 dernieres en 120 facons, qu'on changera en 6 sortes a cause des 6 choses qui doiuent l'vne apres l'autre tenir le 1\textsuperscript{er} lieu.
Pour 7 choses, on fera 720 changem\textsuperscript{ts}. des 6 dernieres, qui seront recommencées 7 fois afin que chacune ayt le 1\textsuperscript{er} lieu et ainsy des aut\textsuperscript{res} a l'infiny.
par ou l'on void qu'il faut commencer à trauailler sur les trois derniers, et gardant la premiere des trois ranger les deux derniers en deux facons, et parce qu'il y a trois lettres, on met chascune des 3 p\textsuperscript{r} la pre\textsuperscript{re}, et apres chascune, les deux autres en deux facons : de la vient que p\textsuperscript{r} auoir la combinaison de 3, on multiplie 2 par 3 p\textsuperscript{r} auoir 6. /
voyez la combinaison des 8 notes de l'octaue dans le gros volume manuscrit entier de ceste matiere, auec les regles des combinaisons deuant et apres. et parce qu'apres auoir consideré la variété de 3 notes en 6 facons, et que chascune de 4 notes peut tenir le 1\textsuperscript{er} rang, on multiplie 6 par 4 p\textsuperscript{r} auoir 24 de combinaison et ainsy de 5, 6, 7 et les autres à l'infiny
\note{« EI » : après « ie change E », la lettre écrite est un « E » suivi d'un « I » ; le passage décrit la permutation de I et E. « le gros volume manuscrit entier de ceste matiere » : renvoi à un autre manuscrit, que la copie ne nomme pas. La page s'achève sur cette phrase ; la vue 21 ouvre une remarque nouvelle, « p. 113. au commencem\textsuperscript{t} ».}

\end{document}
